\documentclass[a4paper]{article}
\usepackage{graphicx}%
\usepackage{multirow}%
\usepackage{amsmath,amssymb,amsfonts}%
\usepackage{amsthm}%
\usepackage{mathrsfs}%
\usepackage[title]{appendix}%
\usepackage{xcolor}%
\usepackage{textcomp}%
\usepackage{manyfoot}%
\usepackage{booktabs}%
\usepackage{algorithm}%
\usepackage{algorithmicx}%
\usepackage{algpseudocode}%
\usepackage{listings}%
%%%%

% Restore itemize indent
\makeatletter
\renewcommand\itemargs{%
   \listfont%
   \leftmargini1.5em%
   \leftmarginii1.25em%
   \leftmarginiii1em%
   \ifnum\@enumdepth=3\topsep0pt\else\ifnum\@enumdepth=2\topsep0pt\else\topsep 6pt\fi\fi%
   \partopsep     \z@%
   \itemsep       \z@%
   \parsep        \z@%
   \labelsep      0.5em%
   \rightmargin   \z@%
   \raggedright%
   \listparindent \z@%
   \itemindent    \z@}%
\makeatother

%%%%%%%%%%%%%%%%%%%%%%%%%%%%%%% Custom packages
\usepackage{array}
\usepackage{amsfonts}
\usepackage{amssymb}
\usepackage{aliascnt}
\usepackage{tikz}
\usetikzlibrary{arrows.meta,positioning}
\usetikzlibrary{calc}
\usepackage{soul}
\usepackage{standalone}
\usepackage{tabularx}
\usepackage{subcaption}

% \usepackage{setspace}
% \doublespacing

% Fix typesetting of vector subscripts in SENN section
\newenvironment{smallpmatrix}
  {\mkern-2mu\left(\mkern-3mu\begin{smallmatrix}}
  {\end{smallmatrix}\mkern-3mu\right)}

% Shorthands for Results tables
\newcommand{\sdag}{\textsuperscript{\dag}}
\newcommand{\sddag}{\textsuperscript{\ddag}}

% Fix table issues with double spacing/referee
\AtBeginEnvironment{longtable}{\def\baselinestretch{1}\normalsize}


% Maths shortcuts
\DeclareMathOperator*{\argmin}{\arg\!\min}
\DeclareMathOperator*{\argmax}{\arg\!\max}
\DeclareMathOperator*{\SGD}{\text{SGD}}
\DeclareMathOperator*{\Opt}{\text{Opt}}
\newcommand\Wnew{\boldsymbol{W}^\textrm{new}}
\newcommand\Wold{\boldsymbol{W}^\textrm{old}}
\newcommand\Wupdated{\boldsymbol{W}^\textrm{upd}}
\newcommand\Hnew{\boldsymbol{h}^\textrm{new}}
\newcommand\Hold{\boldsymbol{h}^\textrm{old}}
\newcommand\Znew{\boldsymbol{z}^\textrm{new}}
\newcommand\Zold{\boldsymbol{z}^\textrm{old}}
\newcommand\A{\mathcal{A}}
\newcommand\T{\mathcal{T}}
\newcommand\D{\mathcal{D}}


% Linear algebra commands (matching Typst conventions)
\newcommand{\trans}[1]{#1^\top}           % Transpose: x^T
\newcommand{\transp}[1]{\left(#1\right)^\top}  % Transpose with parentheses: (x)^T
\newcommand{\pinv}[1]{#1^+}               % Pseudoinverse: x^+
\newcommand{\pinvp}[1]{\left(#1\right)^+} % Pseudoinverse with parentheses: (x)^+
\newcommand{\norm}[1]{\left\lVert#1\right\rVert}
\newcommand{\frob}[1]{\norm{#1}_F} % Frobenius norm with auto-sized delimiters
\newcommand{\sfrob}[1]{\lVert #1 \rVert_F} % Smaller \frob
\newcommand{\trace}[1]{\operatorname{Tr}(#1)} % Trace

% Neuron addition notation
\newcommand{\fanin}{\boldsymbol{\psi}}  % Fan-in weights
\newcommand{\fanout}{\boldsymbol{\omega}} % Fan-out weights
\newcommand{\Fanin}{\boldsymbol{\Psi}}  % Fan-in weights
\newcommand{\Fanout}{\boldsymbol{\Omega}} % Fan-out weights

\newcommand{\postact}[1][]{a^{(#1)}}    % Post-activation
\newcommand{\preact}[1][]{s^{(#1)}}     % Pre-activation
\newcommand{\deltas}{\delta_{s}}   % Pre-activation change
\newcommand{\gradpreact}{\boldsymbol{g}}             % Gradient w.r.t. pre-activation
\newcommand{\Postact}[1][]{\boldsymbol{a}^{(#1)}}    % Batched post-activation
\newcommand{\Preact}[1][]{\boldsymbol{z}^{(#1)}}     % Batched pre-activation
\newcommand{\Gradpreact}{\boldsymbol{G}}             % Batched gradient
\newcommand{\covpostact}[1][]{\boldsymbol{A}^{(#1)}} % Activation covariance
\newcommand{\covgrad}{\boldsymbol{S}}                % Gradient covariance
\newcommand{\Bottleneck}{\boldsymbol{G}^\perp}       % Residual gradient (batched)
\newcommand{\Bmat}[1][]{\boldsymbol{B}^{(#1)}}       % Gradient-activation cross-covariance

\newcommand{\preactext}{s_{\text{ext}}} % Extended pre-activation
\newcommand{\postactext}{a_{\text{ext}}} % Extended post-activation
\newcommand{\Postactext}{\boldsymbol{a_{\text{ext}}}} % Extended batched post-activation
\newcommand{\Preactext}{\boldsymbol{s_{\text{ext}}}} % Extended batched pre-activation
\newcommand{\cext}{C_{\text{ext}}}                   % Number of added neurons

\newcommand{\cp}{C_{l}}                       % Size of the layer after the grown one
\newcommand{\cm}{C_{l-2}}                       % Size of the layer before the grown one
\newcommand{\loss}{\mathcal{L}}                     % Loss function

\newcommand{\isqrtS}{\boldsymbol{S}^{-1/2}}         % Inverse sqrt of S
\newcommand{\sqrtS}{\boldsymbol{S}^{1/2}}           %
\newcommand{\Aext}{\boldsymbol{A}_{\text{ext}}}         % Covariance of extended activations
\newcommand{\isqrtAext}{\Aext^{-1/2}} % Inverse sqrt of A_ext
\newcommand{\sqrtAext}{\Aext^{1/2}}     
\newcommand{\sqrtA}{A^{1/2}}
\newcommand{\scalarp}[1]{\langle #1 \rangle}  % Scalar product

\newcommand{\ie}{\emph{i.e.}}                     % i.e.
\newcommand{\eg}{\emph{e.g.}}

\newcommand{\dtheta}{d\theta}
\newcommand{\EE}{\mathbb{E}}


%% as per the requirement new theorem styles can be included as shown below
\theoremstyle{thmstyleone}%
\newtheorem{theorem}{Theorem}%  meant for continuous numbers
%%\newtheorem{theorem}{Theorem}[section]% meant for sectionwise numbers
%% optional argument [theorem] produces theorem numbering sequence instead of independent numbers for Proposition
\newaliascnt{proposition}{theorem}
\newtheorem{proposition}[proposition]{Proposition}%
\aliascntresetthe{proposition}
\providecommand*{\propositionautorefname}{Proposition}
%%\newtheorem{proposition}{Proposition}% to get separate numbers for theorem and proposition etc.

\theoremstyle{thmstyletwo}%
\newtheorem{example}{Example}%
\newtheorem{remark}{Remark}%

\theoremstyle{thmstylethree}%
\newtheorem{definition}{Definition}%

\raggedbottom
%%\unnumbered% uncomment this for unnumbered level heads




%%%% From the old proof document

\usepackage{amsmath,amssymb,amsthm,mathtools}
\usepackage{longtable}

\numberwithin{equation}{section}
% Keep lemma numbering with theorem while making \autoref print "Lemma".
\newaliascnt{lemma}{theorem}
\newtheorem{lemma}[lemma]{Lemma}
\aliascntresetthe{lemma}
\providecommand*{\lemmaautorefname}{Lemma}
\newaliascnt{corollary}{theorem}
\newtheorem{corollary}[corollary]{Corollary}
\aliascntresetthe{corollary}
\providecommand*{\corollaryautorefname}{Corollary}

\begin{document}

\title[The Neuron Addition Problem: A Survey on Growing Neural Networks]{The Neuron Addition Problem: A Survey on Growing Neural Networks}

%%=============================================================%%
%% GivenName	-> \fnm{Joergen W.}
%% Particle	-> \spfx{van der} -> surname prefix
%% FamilyName	-> \sur{Ploeg}
%% Suffix	-> \sfx{IV}
%% \author*[1,2]{\fnm{Joergen W.} \spfx{van der} \sur{Ploeg} 
%%  \sfx{IV}}\email{iauthor@gmail.com}
%%=============================================================%%
\definecolor{orcidgreen}{HTML}{A6CE39}
\newcommand{\orcidlogo}{%
  \raisebox{-0.1ex}{%
    \begin{tikzpicture}[scale=1,every node/.style={inner sep=0pt,outer sep=0pt}]
      \fill[orcidgreen] (0,0) circle (0.45em);
      \node[white,font=\sffamily\bfseries\tiny] at (0,0) {iD};
    \end{tikzpicture}%
  }%
}
\newcommand{\orcid}[1]{\,\href{https://orcid.org/#1}{\orcidlogo}}


\author*[]{\fnm{Alex} \sur{Davey}$^*$\orcid{0000-0002-2271-4501}}\email{alex.davey@inria.fr}
\equalcont{These authors contributed equally to this work.}
\author*[]{\fnm{Théo} \sur{Rudkiewicz}$^*$\orcid{0009-0003-4374-531X}}\email{theo.rudkiewicz@inria.fr}
\equalcont{These authors contributed equally to this work.}

\author[]{\fnm{Stéphane} \sur{Rivaud}\orcid{0000-0002-1363-442X}}\email{stephane.a.rivaud@inria.fr}
\author[]{\fnm{Stella} \sur{Douka}\orcid{0009-0002-1088-7568}}\email{styliani.douka@inria.fr}
\author[]{\fnm{Félix} \sur{Houdouin}\orcid{0009-0004-4764-410X}}\email{felix.houdouin@polytechnique.edu}
\author[]{\fnm{Julien} \sur{Mille}\orcid{0000-0002-8081-9790}}\email{julien.mille@inria.fr}
\author[]{\fnm{Alena} \sur{Shilova}\orcid{0000-0002-1795-8421}}\email{alena.shilova@inria.fr}
\author[]{\fnm{Guillaume} \sur{Charpiat}\orcid{0009-0003-6000-9410}}\email{guillaume.charpiat@inria.fr}
\author[]{\fnm{Sylvain} \sur{Chevallier}\orcid{0000-0003-3027-8241}}\email{sylvain.chevallier@universite-paris-saclay.fr}

\affil[]{\orgdiv{Université Paris-Saclay}, \orgname{CNRS, INRIA, LISN}, \orgaddress{\city{Orsay}, \postcode{91405}, \country{France}}}

\abstract{
Choosing correct architectures for neural networks is a complex process, and successful architectures often result from great collective efforts, as in vision and language models, or from extensive trial-and-error training. Growing neural network methods propose an alternative: start from a small seed model and expand it during training by \eg{} adding neurons or layers. Commonly cited motivations are (i) reducing training cost compared to full-size training, (ii) better optimization due to the regularisation effect of growth, and (iii) performing a form of architecture search that yields a model that is ``just large enough'' for the task. 
This survey proposes a unified view of a broad class of growing methods, focusing on the problem of neuron addition (how to best choose new neuron parameters). Building on this, we discuss applications to continual learning, reinforcement learning, and transformer growth.
We evaluate current methods on their stated motivations, when and to what degree they are fulfilled, and the future work needed to close existing gaps.
We plan to continue surveying growing methods as a wiki, to which contributions are welcome: \url{https://growingnet.github.io/growing_wiki/}.
}

\keywords{Neural networks, Model structure, Optimization, Energy efficiency}

%%\pacs[JEL Classification]{D8, H51}

%%\pacs[MSC Classification]{35A01, 65L10, 65L12, 65L20, 65L70}

\maketitle

%%%%%%%%%%%%%%%%%%%%%%%%%%%%%%%%%%%%%%%%%%%%%%%%%%%%%%%%%%%%%%%%%%%%%%%%%%%%%%%%%%%%%%%%%%%%%%%%%%%%%%%%%%%%%%%%%%%%%%%%%%%%%%%%%%%%%%%%%%%%%%%%%%%%%%%%%%%%%%%%%%%%%%%%%%%%%%%%%%%%%%%%%%%%%%%%%%%%%%%%%%%%%%%%%%%%%%%%%%%%%%%%%%%%%%%%%%%%%%%%%%%%%%%%%%%%%%%%%%%%%%%%%%%%%%%%%%%%%%%%%%%%%%%%%%%%%%%%%%%%%%%%%%%%%%%%%%%%%%%%%%%%%

% Single author syntax

\section{Introduction}

Neural Networks (NNs) are typically trained by first fixing the architecture $A \in \mathcal{A}$, after which the parameters $\theta \in \Theta_A$ are optimized to minimize a training objective. The architectural choice is crucial, as it defines the function class explored during training, 
and is the workhorse of deep learning~\citep{he_deep_2016,vaswani_attention_2017}. Despite this, architecture search is still typically performed manually, requiring significant expertise and wasted compute due to retraining.

Ideally, we would like to jointly optimize over the architecture space $\mathcal{A}$ and its space of parameters $\Theta_A$ by solving
\begin{align}\label{eqn:ideal_obj}
    A^*, \theta^* = \argmin_{A\in \mathcal{A},\; \theta \in  \Theta_A}  \mathcal{L}(f_{A,\theta}) + \mathcal{C}(A)
\end{align}
where $f_{A,\theta}$ denotes the function induced by architecture $A$ with parameters $\theta$, $\mathcal{L}$ is the empirical risk over a data distribution $\mathcal{D}$, and $\mathcal{C}$ is a penalty on the architecture complexity, measured in total parameters, FLOPs or some other notion of complexity. The closest approach to this objective is \emph{Neural Architecture Search} (NAS)~\citep{zoph_neural_2017}\footnote{NAS does not exactly optimize this objective, since it typically does not find the optimal parameters $\theta^*$.}
as a form of hyper-parameter optimization~\citep{egele:tel-04636586}. However, a full NAS loop is often prohibitively expensive, requiring multiple retrainings, and also ignores a key constraint: we frequently start from a pre-trained model that we would like to adapt rather than discard.

This motivates \emph{growing} Neural Network architectures: starting with a small ``seed'' architecture and expanding its capacity during training by applying local architecture transformations $\T$ (also known as network morphisms), such as widening existing layers or adding new ones, and appropriately adapting the existing weights. Concretely, let $f_{A_t, \theta_t}$ denote the model at growth step $t$ and $\Opt_
\theta(\textrm{goal}(\theta), \theta_{\text{init}})$ denote a few steps of, e.g., stochastic gradient descent to optimize $\text{goal}(\theta)$ over $\theta$ starting from $\theta_{\text{init}}$.
We alternate between applying the growth operator $\T$ and the optimization step
\begin{equation}\label{eqn:grow_decomposition}
\begin{aligned}
    &\theta_t' = \Opt_{\theta}(\mathcal{L}(f_{A_t, \theta}), \theta_t )\\
    &A_{t+1}, \theta_{t+1} = \T(A_t, \theta_{t}') 
\end{aligned}
\end{equation}
in the hope that the final architecture $A_T$ and weights $\theta_T$ are a good approximation to the original objective~\eqref{eqn:ideal_obj}. The growth operator $\T$ is typically constrained to a neighborhood $\mathcal{N}(A_t)$ of the current architecture $A_t$, making local architecture modifications (\eg{}, adding neurons or layers).

In practice, the behavior of $\Opt$ heavily depends on the initialization $\theta_{t+1}$ of the transformed architecture, making it key to achieving good performance. The lottery ticket phenomenon~\citep{frankle_lottery_2019,chen_elastic_2021} highlights that the particular initialization and training path can matter as much as the final architecture, suggesting that growth methods leveraging a fixed set of initial weights have the potential to outperform NAS-like methods, which ignore this.

The abstract formulation of Equation~\eqref{eqn:grow_decomposition} leaves much of the growing problem unspecified, which is often summarized as \emph{where} to grow, \emph{when} to grow, \emph{how} to grow, and \emph{when to stop}. The focus in the literature has been overwhelmingly on \emph{how} to grow, which we term the \emph{neuron addition problem}: when adding new neurons, how many should we add, and how should we initialize them? 

\subsection{Motivations and applications}

Motivations for growing neural networks broadly fall into two categories. In the first, the end point is known: growth is a training strategy for reaching a predefined target architecture $A_T$. Growing is then motivated by faster training~\citep{evci_gradmax_2022,yuan_accelerated_2023}, by improved optimization dynamics~\citep{wu_when_2024,caillon_growing_2021}, or to account for changing data distributions such as in continual learning~\citep{yoon_lifelong_2018,li_learn_2019,yang_grown_2021}.

Alternatively, the endpoint architecture $A_T$ is unknown: growth is used as a frugal form of architecture search, where we would like to discover an architecture that is ``just large enough'' for the task at hand. Such \emph{frugal learning} is becoming increasingly important as achieving state-of-the-art performance increasingly relies on scaling model size and compute, with energy consumption and $\mathrm{CO_2}$ emissions increasing exponentially, outpacing improvements in hardware~\citep{thompson_deep_2021,morand_environmental_2025}. As a frugal method, growing is complementary to compression and pruning, which start from an oversized model and typically only remove parameters after training. While compression and pruning reduce inference costs, they still require training the larger model. The energy efficiency comparison with growing is nuanced, with recent work~\citep{boumendil_grow_2023} showing that the optimal strategy depends on the target accuracy and compute budget, with neither approach dominating across all scenarios.
    
Transformer growth is a prominent application of neural network growth, and is able to significantly reduce the training cost, by growing towards a target full-size architecture~\cite{gu_bert_2021,wang_ligo_2023}. In Sec.~\ref{sec:transformers}, we discuss the connection between transformer growth and the wider growing literature, and empirical evidence for their use.

This survey is \emph{the first modern overview of growing neural networks during training time}, the last one being close to thirty years ago~\citep{kwok_constructive_1997}. In contrast, other surveys target either sparsity and pruning in neural networks~\citep{sparsitySurvey}, inference-time dynamic architectures~\citep{dynamicSurvey}, or comparative studies restricted to transformers~\citep{pandey2024comparative}. Instead, this survey aims to unify the field of growing neural networks, scattered across various growth types, applications, and domains. In particular, we show that several key methods for neuron addition based on function geometry can be unified via a common optimization objective.

This survey is organized as follows. We formalize the neuron addition problem in Sec.~\ref{sec:neuron_addition_problem}, before describing the main classes of neuron initialization, i.e. \emph{how to grow}. We then discuss how the proposed criteria can be used to decide where to grow, when to grow, and when to stop in Sec.~\ref{sec:where_when}. We then broaden the scope beyond width and depth growth to discuss boosting (Sec.~\ref{sec:boosting}), sparse growth and grow-prune methods (Sec.~\ref{sec:sparse_prune_methods}), and then applications to continual learning and reinforcement learning in Sec.~\ref{sec:non_stationary}. Sec.~\ref{sec:results} aggregates a large body of width-growing and depth-growing literature to re-examine the degree to which growing satisfies its stated motivations, including how these conclusions transfer to transformer architecture growth. Finally, we discuss open questions and future directions.


\section{How to grow: The neuron addition problem}\label{sec:neuron_addition_problem}


We consider a feedforward neural network with layers indexed by $l$. The size of layer $l$ is $C_l$, and $n$ denotes the number of data points. The forward pass through consecutive layers is:
\begin{align}
    x \longrightarrow \dots \longrightarrow \postact[l-1](x) \xrightarrow{\;W_l\;} \preact[l](x) \xrightarrow{\;\;\sigma\;\;} \postact[l](x)  \longrightarrow \dots \longrightarrow f(x)
\end{align}
where $\preact[l]: \mathbb{R}^{C_0} \to \mathbb{R}^{C_l}$ and $\postact[l]: \mathbb{R}^{C_0} \to \mathbb{R}^{C_l}$ denote the function giving the pre- and post-activations at layer $l$, respectively. 
For $n$ batched samples, we stack activations row-wise: $\Postact[l] \in \mathbb{R}^{n \times C_l}$. The (negative) gradient of the loss with respect to pre-activations is similarly stacked: $\Gradpreact^{(l)} \in \mathbb{R}^{n \times C_{l}}$.
Without loss of generality, we omit the index $l$ when we consider an object associated with a current layer $l$.

\textbf{Width addition.} We aim to expand layer $l-1$ by adding $k$ new neurons. This is done by adding fan-in weights $\Fanin \in \mathbb{R}^{k \times C_{l-2}}$ and fan-out weights $\Fanout \in \mathbb{R}^{C_l \times k}$:
\begin{align}
    \postact[l-2](x) \overset{\Fanin}{\xrightarrow{\hspace{2em}}} \preactext(x) \overset{\sigma}{\xrightarrow{\hspace{2em}}} \postactext(x) \overset{\Fanout}{\xrightarrow{\hspace{2em}}} \deltas(x)
\end{align}
where $\deltas(x) =  \Fanout \sigma \left(\Fanin \postact[l-2](x) \right)$ and $x$ is the network input.
Hence, the change to the pre-activation at layer $l$ is $\preact[l] \leftarrow \preact[l] + \deltas$, as shown in Fig.~\ref{fig:width-addition}.

\textbf{Layer addition.} Alternatively, we can grow the depth of the network by adding a new layer. In particular, we replace the weights $W_l$ with $W_l \to \Omega \circ \sigma \circ \Psi$, as shown in Fig.~\ref{fig:layer-addition}. Abstractly, one can imagine layer addition as the insertion of a zero-width layer, followed by a width addition step. Although a zero-width layer is not a valid network on its own, this framing provides a natural way to understand how large a new layer to add, and how to initialize these new parameters.

We categorize neuron addition methods by the goal of their initialization strategy:
\begin{figure}
    \centering
    \begin{minipage}[t]{0.48\linewidth}
        \centering
        \includestandalone[width=\linewidth]{figures/neuron_addition_new_r}
        \caption{Width addition in layer $l-1$.}
        \label{fig:width-addition}
    \end{minipage}
    \hfill
    \begin{minipage}[t]{0.48\linewidth}
        \centering
        \includestandalone[width=\linewidth]{figures/layer_addition_new_r}
        \caption{Layer addition in layer $l-1$.}
        \label{fig:layer-addition}
    \end{minipage}
\end{figure}
\begin{itemize}\itemsep0em 
    \item \textbf{Purely function-preserving}: the network output is unchanged after growth (but its gradient may change!).
    \item \textbf{Function geometry-based}: Methods optimizing local objectives such as gradient norm or loss decrease. They can be function-preserving or function-improving.
    \item \textbf{Training dynamics-based}: New weights are initialized to optimize training dynamics.
\end{itemize}

\subsection{Vanilla function-preserving initialization}
\label{sec:function_preserving}

How can we transfer knowledge from one network to a new, larger network? A network morphism $A', \theta' = \T(A, \theta)$ is called \emph{function-preserving} if
\begin{align}
    \forall  x \in \D, \quad f_{A', \theta'}(x) = f_{A, \theta}(x) \, .
\end{align}

\textbf{Net2Net.} In early work, Net2Net~\citep{chen_net2net_2016} proposed a \emph{Net2WiderNet} operation: increasing the width of a layer by splitting the existing weights of a neuron into two new neurons, with the same input weight as the original, but with half the output weights, to compensate for the duplication. In practice, to break symmetry, a small amount of random noise is added to the new neurons.

\textbf{Network Morphism.}~\citet{wei_network_2016} explore other function-preserving network morphisms as an alternative to neuron splitting. Indeed, for a linear network and in the absence of activation functions, any decomposition of the weight matrix $\boldsymbol{W} = \boldsymbol{A}\boldsymbol{B}$ into two shape-compatible matrices is a valid function-preserving morphism. Network Morphism aims to find the least sparse matrix decomposition.

%%%%%%%%%%%%%%%%%%%%%%%%%%%%%%%%%%%%%%%%%%%%%%%%%%%%%%%%%%%%%%%%%%%%%%%%%%%%%%%
\subsection{Exploiting function geometry}
\label{sec:function_geometry}

The methods in this section use local function information to optimize the initialization of new neurons. We define $\Gradpreact \in \mathbb{R}^{n \times \cp}$ as the (negative) gradient of the loss w.r.t.\ pre-activations, stacked over $n$ samples, and the residual gradient $\Bottleneck$ as the gradient component that cannot be addressed by updating existing fan-out weights of the increased layer---it identifies an \emph{expressivity bottleneck}.

The geometrical methods here use either the classical gradient, some form of natural gradient, or both. For classical gradient descent, the gradient is computed for a small, in Euclidean norm, change of parameters $\dtheta$. In contrast, the natural gradient is computed for a small change of the function itself, such that parameters $\theta$ and $\theta'$ are close if $\norm{f_{\theta} - f_{\theta'}}$ is small.  
K-FAC~\citep{dangelKroneckerfactoredApproximateCurvature2025} proposes to approximate this norm when changing $W_l$ to $W_l+\delta W_l$ as $\frob{\sqrtS \delta W_l \sqrtA_{l-1}}$, 
where $A_{l-1} := \postact[l-1] \transp{\postact[l-1]}$ is the covariance of the input of the layer
and $S:= \trans{\frac{\partial f}{\partial \preact[l]}} \frac{\partial f}{\partial \preact[l]}$ is the covariance of the Jacobians of the loss with respect to the layer output.
We can say that $A_{l-1}$ accounts for curvature before layer $l$ and $S$ after.

\subsubsection{Function-preserving geometric methods}

These methods initialize with $\deltas = 0$, preserving the network function, but optimize $(\Fanin, \Fanout)$ to maximize trainability.

\textbf{GradMax}~\citep{evci_gradmax_2022} and SENN~\citep{mitchell_self-expanding_2024} aim to maximise the loss decrease after the next gradient step on the new weights. To achieve this, they maximize the norm of the gradient w.r.t.\ new weights. GradMax initializes $\Fanin = 0$, ensuring $\deltas = 0$, then maximizes the gradient norm w.r.t.\ $\Fanin$ with the constraint $\trans{\Fanout} \Fanout  = I_{\cext}$:
\begin{align}\label{eqn:gradmax}
\Fanout^* = \argmax_{\frob{\Fanout} \leq 1} \frob{\nabla_{\Fanin} \loss(f_{\Fanin, \Fanout})}^2
\end{align}

\textbf{SENN}~\citep{mitchell_self-expanding_2024} extends GradMax by considering a \emph{natural gradient} descent step instead of standard gradient descent and switching from $\Fanin = 0$ to $\Fanout = 0$.
It maximizes, using K-FAC approximation, the gradient for the natural gradient norm of $\begin{smallpmatrix}
    \Fanout \\ W_l
\end{smallpmatrix}$, termed the \emph{natural expansion score}
\begin{equation}\label{eq:senn_main}
\eta_{\text{SENN}} \coloneqq \frob{F^{-\tfrac{1}{2}}\nabla_{\begin{smallpmatrix}
\Fanout{} \\
W
\end{smallpmatrix}}\loss(f)}^2
\end{equation}
where $F$ is (the K-FAC approximation of) the Fisher matrix defined in App.~\ref{sec:senn}. However, in practice, SENN maximizes the norm of the gradient of $\Fanout$ only, backpropagating the residual gradient $\Bottleneck$ to avoid redundancy with existing neurons.

\textbf{NORTH}~\citep{maile_when_2022} takes a geometric perspective: new neurons should add \emph{new directions} in activation space, orthogonal to existing activations. It seeks $\Fanin$ such that new activations $\Postactext$ are orthogonal to existing activations $\Postact[l-1]$.

\subsubsection{Function-improving geometric methods}

These methods initialize with $\deltas \neq 0$, directly improving the loss at initialization.

\textbf{Splitting methods.} One might wonder whether the Net2Net split of one neuron into two, with equally divided weights, is optimal. In S2D~\citep{liu_splitting_2019}, it has been shown that, for an infinitesimal change in the parameters $\|\theta_{t+1} - \theta_t\| \le \epsilon$, this choice of split leads to the fastest decrease of the loss. Consider the post-activation output of a particular neuron $i$ at layer $l-1$. Splitting replaces the original neuron with two neurons:
\begin{align}
    \sigma(\preact[l-1]_i) \to \frac{1}{2} \left( \sigma(\theta_1 \cdot \postact[l-2] ) + \sigma(\theta_2 \cdot \postact[l-2]) \right)
\end{align}
The influence of splitting on the loss is characterized by the minimum eigenvalue $\lambda_{\min}$ of the \emph{splitting matrix}:
\begin{equation}\label{eqn:splitting}
S(\theta) = \underset{x \sim \D}{\mathbb{E}}\left[ \nabla_{\preact[l-1]} \loss(f(x)) \nabla_{\theta\theta}^2 \sigma(\theta \cdot \postact[l-2](x)) \right],
\end{equation}
This ``semi-Hessian'' matrix and its eigenvalues provides a notion of splitting stability: when $\lambda_{\min} > 0$, the loss cannot be improved by splitting; when $\lambda_{\min} < 0$, the maximum decrease is achieved with parameter updates:
\begin{align}
\fanin_{1} &= \theta + \epsilon\, v_{\min}(S(\theta)), \quad
\fanin_{2} = \theta - \epsilon\, v_{\min}(S(\theta))
\end{align}
yielding loss decrease $\Delta \loss \geq \frac{\epsilon^2}{2} \lambda_{\min} + \mathcal{O}(\epsilon^3)$. Since the contribution to loss appears at $\mathcal{O}(\epsilon^2)$, splitting can be thought of as a second-order method to escape local minima. In S2D-E~\citep{wang_energy-aware_2020}, S2D is improved to include energy-aware constraints and a fast gradient-based approximation $S(\theta)$, while S3D~\citep{wu_steepest_2021} generalizes the types of splits considered to non-convex combinations of arbitrary sign.

\textbf{TINY}~\citep{verbockhaven_growing_2024} seeks to find new neurons whose contribution $\deltas$ most directly reduces the loss. Using a first-order Taylor expansion:
\begin{align}
    \mathcal{L}(s + \deltas) = \mathcal{L}(s) + \langle \nabla_s \mathcal{L}, \deltas \rangle + o(\|\deltas\|)
\end{align}
TINY first aligns the contribution of existing parameters with $\nabla_s \mathcal{L}$, then the contribution of the new neurons with the residual gradient $\Bottleneck$ to avoid redundancy with existing neurons. Linearizing the activation function $\sigma(x)$ around $0$, this becomes a low-rank matrix approximation:
\begin{align}\label{eq:tiny_main}
\Fanin^*, \Fanout^* = \argmin_{\Fanin, \Fanout} \frob{\Bottleneck - \Postact[l-2] \trans{\Fanin} \trans{\Fanout}}^2
\end{align}
solved in closed form using two SVDs.

\textbf{NeST}~\citep{daiNeSTNeuralNetwork2019} targets sparse networks where connections are added incrementally. Justified by Hebbian theory ``Neurons that fire together wire together''~\citep{lowel_1992}, NeST looks at the correlation between the activation of neuron $i$ in layer $l-2$ and the gradient of neuron $j$ in layer $l$ to decide which neuron to add. They add a neuron with a single non-zero fan-in weight and a single non-zero fan-out weight.
NeST uses the same mechanism to densify layers except that they consider interaction between two successive layers, \eg{} $l-1$ and $l$.

\textbf{Firefly}~\citep{wu_firefly_2020} combines splitting existing neurons, adding new neurons, and adding new layers. It optimizes the candidate parameters through classical gradient descent and estimates the contribution of each candidate neuron. Finally, only the architecture changes with the largest contributions are kept.


%%%%%%%%%%%%%%%%%%%%%%%%%%%%%%%%%%%%%%%%%%%%%%%%%%%%%%%%%%%%%%%%%%%%%%%%%%%%%%%
\subsubsection{A unified framework for geometric methods}
\label{sec:unified}

Many of the above methods, even if they seem different at first sight, can be unified under the same objective with different choices of target, metric, and constraints. Following the work of \citep{verbockhavenSpottingExpressivityBottlenecks2025} that compares TINY to GradMax and NORTH, we propose the following theorem. Proofs can be found in App.~\ref{app:proofs}.
\begin{theorem}[Unified objective for neuron addition]\label{thm:unified_obj_main_text}
Many of the gradient-informed methods above can be derived from the following optimization problem:
\begin{align}\label{eqn:unified_obj}
& \argmin_{\Fanin, \Fanout} \frob{T \isqrtS - \sigma(\Postact[l-2] \trans{\Fanin}) \trans{\Fanout} \sqrtS}\\
&\propto \argmax_{\Fanin, \Fanout; \frob{\sqrtAext \trans{\Fanout} \sqrtS} \le 1} \scalarp{T \isqrtS, \frac{1}{n} \Postactext \trans{\Fanout} \sqrtS}
\end{align}
where:
\begin{itemize}
    \item $T \in \mathbb{R}^{n \times \cp}$ is the target: gradient $\Gradpreact$ or residual gradient $\Bottleneck$
    \item $\covgrad \in \mathbb{R}^{\cp \times \cp}$ is the post-network metric $\covgrad := \EE_x\left[\transp{\frac{\partial f}{\partial \preact[l]}(x)} \frac{\partial f}{\partial \preact[l]}(x) \right]$
    \item $\Aext \in \mathbb{R}^{\cext \times \cext}$ is the activation covariance of the extended activations $\Aext:= \frac{1}{n} \transp{\Postactext} \Postactext$ and $\widehat{\Aext}$ is its approximation for a linearized activation function (which is justified for $\Fanin \approx 0$)
    \item $\sigma$ is the activation function (or its linear approximation)
    \item Additional constraints may be imposed on $\Fanin$ and $\Fanout$.
\end{itemize}

\noindent Table~\ref{tab:unified} summarizes the specialization for each method. We also add a column indicating whether the method tries to satisfy the NORTH orthogonality constraint $\Postactext \perp \Postact[l-1]$ (in the case of TINY it is only satisfied for linearized activation functions).
\end{theorem}

\begin{table}[t]
\centering
\caption{Neuron addition methods within the unified objective~\eqref{eqn:unified_obj}.}
\label{tab:unified}
\small
\begin{tabular}{lccccccc}
\toprule
\textbf{Method} & $T$ & $\covgrad$ & $\Aext$ & \textbf{Constraint} & $\Fanin$ & $\Fanout$ & $\Postactext \perp \Postact[l-1]$ \\
\midrule
TINY  & $\Bottleneck$ & $I_{\cp}$ & $\widehat{\Aext}$ & -- & $\Fanin^*$ & $\Fanout^*$ & Yes \\
GradMax & $\Gradpreact$ & $I_{\cp}$ & $\Fanin \trans{\Fanin}$ & $\trans{\Fanout}\Fanout = I$ & $0$ & $\Fanout^*$ & No \\
SENN & $\Bottleneck$ & $\covgrad$ & $\Aext$ & -- & $\Fanin^*$ & $0$ & Yes \\
NeST & $\Gradpreact$ & $I_{\cp}$ &  $\Fanin \trans{\Fanin}$ & $\|\Fanin\|_0 = \|\Fanout\|_0 = 1$ & $\Fanin^*$ & $\Fanout^*$ & No \\
\bottomrule
\end{tabular}
\vspace{-1em}
\end{table}

The key insight is that all methods minimize the distance between the new neurons' contribution and some target gradient signal, differing in:
\begin{enumerate}
    \item \textbf{Target}: whether they use the raw gradient $\Gradpreact$ or residual gradient $\Bottleneck$ (which avoids redundancy with existing neurons).
    \item \textbf{Metric} on parameters: whether they use natural norm integrating pre- and post-network curvature, semi-natural norm only integrating pre-network curvature ($\covgrad = I_{\cp}$), or standard Euclidean norm
     ($\covgrad = I_{\cp}$ and $\Aext = I_{\cext}$).
    \item \textbf{Constraints}: sparsity ($\|\cdot\|_0 = 1$), or function preservation ($\Fanin = 0$ or $\Fanout = 0$).
\end{enumerate}

%%%%%%%%%%%%%%%%%%%%%%%%%%%%%%%%%%%%%%%%%%%%%%%%%%%%%%%%%%%%%%%%%%%%%%%%%%%%%%%
\subsection{Improving training dynamics}
\label{sec:training_dynamics}

Instead of focusing solely on function preservation, or immediate improvements in e.g. the loss, Variance Transfer~\citep{yuan_accelerated_2023} uses a function-preserving initialization with the weight distribution and learning rate chosen to promote future training dynamics. Note that, for any matrices $V \in \mathbb{R}^{k/2\times C_{l-2}}$ and $Z \in \mathbb{R}^{C_l \times k/2}$, the addition of new neurons with a minus sign inserted
\begin{align}
    \Fanin = \begin{bmatrix} V \\ V \end{bmatrix}, \qquad
    \Fanout = \begin{bmatrix} Z & -Z \end{bmatrix}
\end{align}
ensures that the contributions of the new weights cancel, preserving the network function. Inspired by studies on how to optimally transfer hyperparameters across networks of different widths~\citep{yang_tensor_2021}, Variance Transfer initializes the weights of the final output layer as $Z \sim \mathcal{N}(0, 1/\textrm{fan\_in}^2)$ rather than the usual $1/\textrm{fan\_in}$. When the weight matrix is enlarged, layer weights are rescaled by $\boldsymbol{W}_{t+1}=\boldsymbol{W}_t \cdot \frac{C_t}{C_{t+1}}$ to preserve variance.

Furthermore, the learning rate is adapted to the growth cycle of each sub-network. Partitioning the weights $\mathbf{W}_T$ of the entire network according to the growth stage $t \in [0, T]$ at which they were added, $\mathbf{W} = \{W_0, W_{\Delta 1}, \ldots, W_{\Delta T} \}$, each sub-network $W_{\Delta t}$ is assigned a learning rate:
\begin{align}
\eta_t = \eta_0 \, \frac{\|\boldsymbol{W}_{\Delta t}\|_F}{\|\boldsymbol{W}_{0}\|_F}
\end{align}
where $\eta_0$ is the base learning rate and $\|\cdot\|_F$ is the Frobenius norm. This compensates for the fact that different subnetworks are trained for a different number of epochs. These modifications to the vanilla function-preserving morphism result in a surprisingly strong baseline.

Tables~\ref{tab:growing_methods} and~\ref{tab:growing_decisions} provide a summary of the key growth metric for each of the main width and depth growing methods, as well as how each one answers the key questions of: how to grow, where to grow, when to grow, and when to stop.

\newcommand{\growingdash}{\multicolumn{1}{c}{--}}

\begin{table*}[t!]
\centering
\caption{Comparison of growing methods by their principal metric or objective, growth mechanism, and supported type of architectural growth.}
\label{tab:growing_methods}
\fontsize{5.5pt}{6.5pt}\selectfont
\setlength{\tabcolsep}{3pt}
\begin{tabularx}{\linewidth}{
    p{0.20\linewidth}
    >{\centering\arraybackslash}p{0.14\linewidth}
    >{\centering\arraybackslash}p{0.16\linewidth}
    >{\centering\arraybackslash}p{0.07\linewidth}
    >{\centering\arraybackslash}X
    >{\centering\arraybackslash}p{0.07\linewidth}
    >{\centering\arraybackslash}p{0.07\linewidth}
}
\toprule
\textbf{Method}
& \multicolumn{3}{c}{\textbf{Growth Metric}}
& \textbf{How}
& \textbf{Width}
& \textbf{Depth} \\

\midrule

\textbf{Net2Net}
& \growingdash
& \growingdash
& \growingdash
& \mbox{function preserving}
& $\checkmark$
& $\checkmark$ \\

\textbf{Net. Morph.}
& \growingdash
& \growingdash
& \growingdash
& \mbox{function preserving}
& $\checkmark$
& $\checkmark$ \\

\textbf{Variance Transfer}
& \growingdash
& \growingdash
& \growingdash
& \mbox{function preserving}
& $\checkmark$
& \\

\textbf{AdaNet}
& \mbox{risk--complexity}
& $F(\mathbf{w})$
& \eqref{eq:adanet_main}
& $\min F(\mathbf{w})$
& $\checkmark$
& $\checkmark$ \\

\midrule

\textbf{GradMax}
& \mbox{gradient norm}
& $\|\nabla_{\Fanin}\loss(f)\|$
& \eqref{eqn:gradmax}
& $\max\|\nabla_{\Fanin}\loss(f)\|$
& $\checkmark$
& \\

\textbf{TINY}
& \mbox{exp. bottleneck}
& $\|\Bottleneck\|$
& \eqref{eq:tiny_main}
& $\operatorname{SVD}(\Bottleneck)$
& $\checkmark$
& \\

\textbf{NeST}
& \mbox{act.--grad. correlation}
& $|B_{-2}|$
& \growingdash
& \mbox{sparse fan-in/out init.}
& $\checkmark$
& \\

\textbf{SENN}
& \mbox{expansion score}
& $\Delta\eta_{\mathrm{SENN}}$
& \eqref{eq:senn_main}
& $\max\Delta\eta_{\mathrm{SENN}}$
& $\checkmark$
& $\checkmark$ \\

\textbf{Splitting}
& \mbox{min. eigenvalue}
& $\lambda_{\min}(S(\theta))$
& \eqref{eqn:splitting}
& $\mathbf v_{\min}(S(\theta))$
& $\checkmark$
& \\

\textbf{Firefly}
& \mbox{neuron contribution}
& $s_i$
& \growingdash
& $\max|s_i|$
& $\checkmark$
& $\checkmark$ \\

\textbf{NORTH}
& \mbox{eff. rank}
& $\phi_a^{\mathrm{ED}}$
& \eqref{eq:north_main}
& $\Postactext\perp\Postact[l-1]$
& $\checkmark$
& \\

\midrule

\textbf{LipGrow}
& \mbox{scaled Lipschitz}
& $L(F)/L_0$
& \growingdash
& \mbox{clone weights}
&
& $\checkmark$ \\

\textbf{FRAGrow}
& \mbox{overfitting risk}
& $\mathrm{ORL}$
& \eqref{eq:fragrow_main}
& \mbox{clone weights}
&
& $\checkmark$ \\

\textbf{AutoGrow}
& \mbox{validation gain}
& $\Delta\mathrm{Acc}_{\mathrm{val}}$
& \growingdash
& gaussian
&
& $\checkmark$ \\

\textbf{SensLI}
& \mbox{layer sensitivity}
& $\|\nabla_{\theta^\mathrm{new}_l}\mathcal L\|$
& \eqref{eq:sensli_main}
& $f'=f$, $\nabla_{\theta_{\rm new}}\mathcal L\ne0$
&
& $\checkmark$ \\

\bottomrule
\end{tabularx}
\end{table*}

\begin{table*}[b!]
\centering
\caption{Summary of the criteria used by growing methods to decide when to grow, where to grow and when to stop. Here, $\tau_g$, $\tau_\ell$, and $\tau_s$ denote generic method-specific thresholds for growth, location eligibility, and stopping, respectively; identical subscripts indicate the same threshold. A dash denotes non-adaptive behavior: a predefined schedule for \emph{when to grow} and \emph{when to stop}, and growing either everywhere at once or in a round-robin fashion for \emph{where to grow}.}
\label{tab:growing_decisions}
\fontsize{5.5pt}{6.5pt}\selectfont
\setlength{\tabcolsep}{3pt}
\begin{tabularx}{\linewidth}{
    p{0.20\linewidth}
    >{\centering\arraybackslash}p{0.14\linewidth}
    >{\centering\arraybackslash}p{0.14\linewidth}
    >{\centering\arraybackslash}p{0.05\linewidth}
    >{\centering\arraybackslash}p{0.14\linewidth}
    >{\centering\arraybackslash}X
    >{\centering\arraybackslash}p{0.12\linewidth}
}
\toprule
\textbf{Method}
& \multicolumn{3}{c}{\textbf{Growth Metric}}
& \textbf{When}
& \textbf{Where}
& \textbf{Stop} \\
\midrule

\textbf{Net2Net}
& \growingdash & \growingdash & \growingdash & \growingdash & \growingdash & \growingdash \\

\textbf{Net. Morph.}
& \growingdash & \growingdash & \growingdash & \growingdash & \growingdash & \growingdash \\

\textbf{Variance Transfer}
& \growingdash & \growingdash & \growingdash & \growingdash & \growingdash & \growingdash \\

\textbf{AdaNet}
& \mbox{risk--complexity}
& $F(\mathbf{w})$
& \eqref{eq:adanet_main}
& $F(\mathbf{w})<\tau_g$
& $\min F(\mathbf{w})$
& $F(\mathbf{w})\ge\tau_g$ \\

\midrule

\textbf{GradMax}
& \mbox{gradient norm}
& $\|\nabla_{\Fanin}\loss(f)\|$
& \eqref{eqn:gradmax}
& \growingdash & \growingdash & \growingdash \\

\textbf{TINY}
& \mbox{exp. bottleneck}
& $\|\Bottleneck\|$
& \eqref{eq:tiny_main}
& \growingdash & \growingdash & \growingdash \\

\textbf{NeST}
& \mbox{act.--grad. correlation}
& $|B_{-2}|$
& \growingdash
& \growingdash
& $\max |B_{-2}|$
& \growingdash \\

\textbf{SENN}
& \mbox{expansion score}
& $\Delta\eta_{\mathrm{SENN}}$
& \eqref{eq:senn_main}
& $\Delta\eta_{\mathrm{SENN}}>\tau_g$
& $\max\Delta\eta_{\mathrm{SENN}}$
& $\Delta\eta_{\mathrm{SENN}}\le\tau_g$ \\

\textbf{Splitting}
& \mbox{min. eigenvalue}
& $\lambda_{\min}(S(\theta))$
& \eqref{eqn:splitting}
& $\lambda_{\min}<0$
& $\min\lambda_{\min}$
& $\lambda_{\min}\ge0$ \\

\textbf{Firefly}
& \mbox{neuron contribution}
& $s_i$
& \growingdash
& \growingdash
& $\max |s_i|$
& \growingdash \\

\textbf{NORTH}
& \mbox{eff. rank}
& $\phi_a^{\mathrm{ED}}$
& \eqref{eq:north_main}
& $\phi_a^{\mathrm{ED}}>\tau_g$
& $\phi_a^{\mathrm{ED}}>\tau_\ell$
& \growingdash \\

\midrule

\textbf{LipGrow}
& \mbox{scaled Lipschitz}
& $L(F)/L_0$
& \growingdash
& $L(F)/L_0>\tau_g$ & \growingdash & \growingdash \\

\textbf{FRAGrow}
& \mbox{overfitting risk}
& $\mathrm{ORL}$
& \eqref{eq:fragrow_main}
& $I(\mathrm{ORL})$ & \growingdash & \growingdash \\

\textbf{AutoGrow}
& \mbox{validation gain}
& $\Delta\mathrm{Acc}_{\mathrm{val}}$
& \growingdash
& \growingdash
& \growingdash
& $\Delta\mathrm{Acc}_{\mathrm{val}}<\tau_s$ \\

\textbf{SensLI}
& \mbox{gradient norm}
& $\|\nabla_{\theta^\mathrm{new}_l}\mathcal L\|$
& \eqref{eq:sensli_main}
& $\|\nabla_{\theta^\mathrm{new}_l}\mathcal L\|\ge\tau_g$
& $\max \|\nabla_{\theta^\mathrm{new}_l}\mathcal L\|$
& \growingdash \\

\bottomrule
\end{tabularx}
\end{table*}


%%%%%%%%%%%%%%%%%%%%%%%%%%%%%%%%%%%%%%%%%%%%%%%%%%%%%%%%%%%%%%%%%%%%%%%%%%%%%%%
\subsection{Depth growth}
\label{sec:adding_layers}

Depth is grown by three operators, treated in turn. Serial deepening replaces a layer on the main path by a short sequence. Residual insertion replaces an activation by a shape-preserving branch. Parameter reuse builds the new depth from weights that have already been trained.

\subsubsection{Serial deepening}

The most direct way to add depth is to replace a single layer by a short sequence of layers on the same path. Following Fig.~\ref{fig:layer-addition}, the affine map $u\mapsto W_lu$, evaluated at $u=\postact[l-1](x)$, becomes
\begin{align}
u\mapsto \Omega\,\sigma(\Psi u).
\end{align}
All the methods below keep the trained map in the first factor, $\Psi=W_l$, and differ only in their choice of the new matrix $\Omega$. Preserving the function constrains that choice through the product $\Omega\Psi=W_l$, a constraint the methods below either enforce or deliberately break.

\textbf{Net2DeeperNet}~\citep{chen_net2net_2016} uses the identity fill $(\Psi,\Omega)=(W_l,I)$, so that the inserted path reduces to $u\mapsto\sigma(W_lu)$. When the original layer already applies $\sigma$ after $W_lu$, and the activation is idempotent, $\sigma\circ\sigma=\sigma$, as is the case for ReLU, this post-activation is left unchanged and the morphism is function-preserving.

\textbf{Network Morphism}~\citep{wei_network_2016} relaxes the identity fill to any factorization satisfying $\Omega\Psi=W_l$ (or its convolutional analogue), thereby allowing a denser pair than $(W_l,I)$. When $\sigma$ is not idempotent, they replace it by the parametric family $\sigma^a=(1-a)\sigma+a\,\mathrm{id}$, initialized at $a=1$ so that $u\mapsto\Omega\sigma^a(\Psi u)$ again reduces to $W_lu$ at insertion, and then train $a$ with the task loss.

The same serial topology can also be used with an $\Omega$ that does not satisfy $\Omega\Psi=W_l$, in which case the function is not preserved. On PlainNet architectures, AutoGrow~\citep{wen_autogrow_2020} uses a random fill $(\Psi,\Omega)=(W_l,\Omega_{\mathrm{rand}})$, whereas FRAGrow~\citep{wu_when_2024} copies the trained layer, $(\Psi,\Omega)=(W_l,W_l)$. Such constructions remain the operator of choice for plain, non-residual networks; in modern residual architectures they are largely superseded by the insertions described next.

\subsubsection{Residual insertion}

Given activation $h$ of a layer, we can replace it by a pair of activations $h_{\mathrm{in}}$ and $h_{\mathrm{out}}$ such that
\begin{align}\label{eqn:residual_insertion}
h_{\mathrm{out}} = h_{\mathrm{in}} + F(h_{\mathrm{in}}),\qquad h_{\mathrm{in}}=h,
\end{align}
Compared with serial deepening, $F$ need only be shape-preserving: its input and output dimensions must match, which excludes downsampling blocks. Furthermore, function-preservation only requires $F(h_{\mathrm{in}}) = 0$.

\textbf{Inheriting from width growth.} When $F$ is a single hidden layer, any width-expansion method that supports growth from an empty hidden layer canonically becomes a depth-growth method via residual insertion. The unified geometrical view of first-order width-growth methods is inherited as well. Theorem~\ref{thm:unified_obj_main_text} applies with the incoming activations $\Postact[l-2]$ read as $h_{\mathrm{in}}$, when the skip connection is the identity. This recovers the width-addition update of Sec.~\ref{sec:neuron_addition_problem},
\begin{align}\label{eqn:residual_width}
\preact[l] \leftarrow \preact[l] + \deltas, \qquad \deltas = \Fanout\sigma(\Fanin \postact[l-2]).
\end{align}
An empty layer, taken in isolation, has an empty tangent space, so $\Bottleneck=\Gradpreact$ for TINY~\citep{verbockhaven_growing_2024} and the orthogonality constraint of NORTH~\citep{maile_when_2022} is vacuous. GradMax~\citep{evci_gradmax_2022} takes $\Fanin=0$ and optimizes $\Fanout$; SENN~\citep{mitchell_self-expanding_2024} and GradMax-Opt take $\Fanout=0$ and optimize $\Fanin$.

\textbf{Function preservation.} Initializing $F(h_{\mathrm{in}})=0$ is enough to preserve the function, but a gradient must still be able to flow through the new parameters if the expansion is to be trainable. For a single hidden layer one may set $\Fanin=0$ or $\Fanout=0$, but not both, or both gradients are zero and remain so.\footnote{$\Fanout = 0$ preserves the function for any activation, whereas $\Fanin = 0$ also requires $\sigma(0) = 0$ and a zero branch bias, which holds for ReLU, GELU and tanh but not for the sigmoid.} The same one-sided silence is often realized by zeroing the last linear map, or the affine parameters of the last normalization layer, so that earlier weights can remain nonzero and carry gradient: SensLI~\citep{kreisSensLISensitivityBasedLayer2026} zeros $W_2$; AutoGrow~\citep{wen_autogrow_2020} the BatchNorm affine $\gamma,\beta$; Gesmundo--Maile~\citep{gesmundo_composable_2023} zeros the MHA and MLP projections; and $G_{\mathrm{zero}}$~\citep{du_stacking_2024} zeros the new output parameters.

\textbf{Masked growth.} Gating the insertion as $h_{\mathrm{out}}=h_{\mathrm{in}}+\alpha F(h_{\mathrm{in}})$ with $\alpha=0$ leaves $F$ unconstrained. MSG~\citep{yao_msg_2024} grows transformers this way, raising a transient mask linearly to $1$ over a fixed number of steps.

\textbf{Neural ODE.} A residual block $h_{\mathrm{out}}=h_{\mathrm{in}}+F(h_{\mathrm{in}})$ can be viewed as one Euler step of $\dot{h}=F(h)$. LipGrow~\citep{dong_towards_2020} replaces that step by two half-steps (clone $F$, rescale the increment), i.e.\ a finer discretization of the same continuous flow. That is only approximately function-preserving, and only while $\mathrm{Lip}(F)$ stays small---which is why they use the Lipschitz constant as a when-to-grow trigger (Sec.~\ref{sec:where_when}).

\subsubsection{Reusing parameters}

The last family builds the new depth from parameters that have already been trained. Unlike residual insertion, we do not have $F(h_{\mathrm{in}})=0$. Let $B=(B_1,\ldots,B_L)$ be a sequence of blocks with parameters $\theta=(\theta_1,\ldots,\theta_L)$, and write $f_B=B_L\circ\cdots\circ B_1$ for the function output. A reuse operator produces a longer sequence $B'=(B'_1,\ldots,B'_{L'})$, $L'>L$, by a linear mix
\begin{align}\label{eqn:reuse}
\theta'_k = \sum_{j=1}^{L} w_{kj}\,\theta_j, \qquad k=1,\ldots,L'.
\end{align}
The grown blocks sit at depths for which those weights were not trained, so $f$ changes. The exception is $G_{\mathrm{zero}}$~\citep{du_stacking_2024}, already discussed above: it is residual silence, not reuse. Copy and nearest-neighbor interpolation are the 0-1 special cases $w_{k,\pi(k)}=1$, i.e.\ $\theta'_k=\theta_{\pi(k)}$; a learned $w$ is the remainder of the family.

\textbf{Copy.} When $w$ is a 0-1 assignment, each new block is a copy of a source block. The default is to compose the trained stack with itself,
\begin{align}\label{eqn:gstack}
f_{B'} = f_B\circ f_B, \qquad \pi=(1,\ldots,L,1,\ldots,L),
\end{align}
which is $G_{\mathrm{stack}}$~\citep{du_stacking_2024} and the depth operator of Progressive BERT~\citep{gu_bert_2021}, usually applied in stages: train $L$ blocks, stack, continue. Because $f_B\circ f_B\neq f_B$, the grown network is not $f$. The remaining choices are which window to copy and where to put it. MIDAS~\citep{saunshi_inductive_2024} takes $\pi$ to repeat a contiguous middle segment and expands from the middle rather than appending on top, which they motivate as unrolling a looped block. Copies inherit the shapes of the source blocks, so the residual-stream dimension is left unchanged.

\textbf{Interpolate.} Interpolation uses the same 0-1 rule with a different $\pi$: new blocks are inserted \emph{between} existing ones and each is filled from its nearest trained neighbor, rather than appending a copy of a window. AutoProg~\citep{li_autoprog_2022} is this nearest-neighbor $\pi$, implemented as a weight-nesting elastic supernet: a subnetwork runs a subset of the intercalated blocks, so growing depth is selecting a denser subset rather than allocating a new tensor. The source need not be the live weights. MoGrow~\citep{li_autoprog_2022} applies the same $\pi$ to an exponential moving average of $\theta$, which reduces the jump in $f$ without forcing the new blocks to be silent.

\textbf{Learned map.} Copy and nearest-neighbor interpolation fix $w$ in advance. LiGO~\citep{wang_ligo_2023} learns it, so~\eqref{eqn:reuse} is $\theta'=M(\theta)$ with $M$ the linear map given by $w$. A dense $w$ would couple every source block to every target block; they factor $M$ as a depth mixture $w\otimes I$ and an optional width Kronecker map, which makes $M$ cheap enough to train. The 0-1 assignments above are recovered by particular $w$. $M$ is fit by a short inner loop on the training loss of the grown network and then discarded: $\theta'$ is only an initialization. The $G_{\mathrm{learn}}$ operator of \citet{du_stacking_2024} is this family.

Quantitative comparisons of depth-growing methods are collected in Table~\ref{tab:results_depth_growing}, while the timing and placement of growth events are treated in Sec.~\ref{sec:where_when}, and their transformer-specific realizations in Sec.~\ref{sec:transformers}. Growing a topology that is not a chain requires an additional choice of location, which we discuss in Sec.~\ref{sec:dag_growth}.

\subsection{Growing Transformers}\label{sec:transformers}

Transformer growth follows essentially the same growing approaches as the wider literature, described above, but one can grow along multiple dimensions. Width growth may increase the token embedding dimension, the intermediate FFN dimension, or the number of self-attention heads, while depth growth adds complete transformer blocks~\citep{yao_msg_2024}. Growth can also be combined with increases in sequence length or image resolution~\citep{gu_bert_2021,li_autoprog_2022}.

For depth growth, a common approach is to initialize added blocks by copying trained blocks, often termed \emph{stacking}~\citep{gong_stacking_2019,du_stacking_2024}. Progressive Stacking repeatedly duplicates the network to increase its depth~\citep{gong_stacking_2019}, while MIDAS grows by copying subsets of blocks and studies insertion within the network~\citep{saunshi_inductive_2024}. For width growth, bert2BERT adapts Net2Net's function-preserving parameter duplication and additionally uses information from adjacent layers; it combines widening with depth expansion by stacking~\citep{chen_bert2bert_2022}.
LiGO instead learns a linear parameter-transfer map, factorized across width and depth using Kronecker structure~\citep{wang_ligo_2023}.
These approaches address the initialization question of Sec.~\ref{sec:neuron_addition_problem} through direct reuse or learned combinations of existing weights.

Several methods are purely function-preserving. LEMON initializes new parameters to ensure function-preservation, accounting for normalisation when widening and using zero-initialized residual contributions when deepening~\citep{wang_lemon_2023}.
MSG uses masks to initially suppress the contribution of new components and progressively activates them during training~\citep{yao_msg_2024}, ensuring function preservation regardless of the choice of weight initialization.

Staged Training~\citep{shen_staged_2022} treats growth as a transformation of the training state: it transfers optimizer statistics alongside model weights and adjusts the learning-rate schedule to support continued training. In the same spirit as Variance Transfer (Sec.~\ref{sec:training_dynamics}),
SPARKLING focuses on preserving hidden-activation scales, expressed through their root mean square, and breaking the symmetry between duplicated parameters~\citep{yu_sparkling_2026}.
It retains the optimizer state for existing weights, and resets it only for new weights so that the duplicates receive different optimizer updates.


\section{Where to grow, when to grow, and when to stop}\label{sec:where_when}


\subsection{Where to grow?}
The Where column of Table~\ref{tab:growing_decisions} summarizes the available location criteria. Most methods do not propose a novel \emph{where to grow} criterion, and either grow everywhere at once~\citep{chen_net2net_2016,wei_network_2016,yuan_accelerated_2023,evci_gradmax_2022,dong_towards_2020} or in a predefined order such as round-robin~\citep{verbockhaven_growing_2024,wen_autogrow_2020,wu_when_2024}. Beyond these baselines, two families of methods stand out. The first reuses the local expressivity criteria already used to decide
when to grow (and for some, when to stop): SENN~\citep{mitchell_self-expanding_2024}
grows where the natural expansion increment $\Delta\eta_{\mathrm{SENN}}$ is largest,
 NORTH~\citep{maile_when_2022} where its activation effective rank or gradient trigger
exceeds a threshold, Splitting~\citep{liu_splitting_2019} splits the
neurons whose splitting matrix $\mathcal{S}(\theta)$ has a the most negative minimum eigenvalue $\lambda_{\min}$, and
\citet{douka_growth_2025} extend TINY's expressivity bottleneck to architectures
with an arbitrary directed acyclic graph structure and grows at its maximum location. The second ranks locations
by measuring the norm of the gradient of new parameters: GradMax~\citep{evci_gradmax_2022} proposes,
without implementing it, to use the singular values returned by its
own maximisation problem. SensLI~\citep{kreisSensLISensitivityBasedLayer2026} instead uses the relative layer sensitivity
\begin{equation}\label{eq:sensli_main}
\rho_{\mathrm{SensLI}}(l)
:=
\frac{\|\nabla_{\theta_\mathrm{new}^l}\mathcal L\|^2}
{|\mathcal W|^{-1}\sum_{W\in\mathcal W}\|\nabla_W\mathcal L\|^2},
\end{equation}
where $\mathcal W$ is the set of existing weight tensors, and inserts a function-preserving layer at the position where this sensitivity is largest.

\subsection{When to grow?}

Table~\ref{tab:growing_decisions} gives a compact overview of the timing and termination criteria for key growing methods. Almost all methods interleave growth events with conventional training, so the interval between two consecutive events is a design choice in its own right. What constitutes a good answer depends on the purpose assigned to growth: either reaching a \emph{predefined target architecture} at a lower training cost, or with a better final performance, or \emph{discovering} an architecture suited to a task for which no target is known. Most methods simply grow on a predefined schedule~\citep{chen_net2net_2016,wei_network_2016,evci_gradmax_2022,verbockhaven_growing_2024,yuan_accelerated_2023}. However, the depth-growing literature in particular focuses on \emph{when to grow}.

\subsubsection{Predefined target architecture}

When the final architecture is specified by the user, growth is only a means of
training it, and the question of when to grow reduces to a pacing problem: how
quickly should a known growth budget be spent?

\textbf{LipGrow}~\citep{dong_towards_2020} grows by duplicating residual blocks and triggers growth when its scaled Lipschitz constant exceeds a tolerance: $L(F)/L_0 >r_{\mathrm{tol}}$, where $L_0$ is the pre-growth Lipschitz constant. An alternative reading of the same criterion, not the one advanced by the
authors, is that a rapidly increasing Lipschitz constant indicates that the network is beginning to overfit the current incompletely-grown architecture, and that growth is preferable before the current weights settle in such a region.

\textbf{FRAGrow}~\citep{wu_when_2024} is explicitly designed to regulate overfitting. Growth acts as an implicit regularizer: blocks added early have time to overfit more, compared to blocks added later. FRAGrow measures the overfitting risk level and uses it to set the growth interval:
\begin{align}\label{eq:fragrow_main}
\mathrm{ORL}
&:=\mathrm{Acc}_{\mathrm{train}}-\mathrm{Acc}_{\mathrm{val}},
& I(\mathrm{ORL})
&:=\frac{I_{\max}}{1+\exp(\alpha-\mathrm{ORL})}.
\end{align}
It grows quickly on models that underfit and slowly on models that overfit. In its primary implementation, each new block copies the weights of its preceding block; only a block following a downsampling block is randomly initialized.

\subsubsection{Unknown target architecture}

When the architecture is an output of the procedure rather than an input, capacity should be added only when it is
needed. The most direct implementation of this idea is to train until
convergence, grow, and retain the change if it improves performance.

\textbf{AutoGrow} \citep{wen_autogrow_2020} (and FRAGrow \citep{wu_when_2024})
report that waiting for convergence is suboptimal, which is consistent with the
observation that converged weights are a poor warm start for further training
\citep{ash_warm-starting_2020}. AutoGrow therefore grows every $K$ epochs,
with $K$ small enough that growth precedes convergence of the current
sub-network. The resulting networks are well trained, in the
sense that their accuracy matches or exceeds that of the same architectures trained from
scratch, but because blocks are added before any expressivity deficiency has been
established, the discovered architectures tend to be much larger than necessary.

The alternative is to detect expressivity deficiencies locally and to grow when,
and where, they arise. \emph{NORTH} \citep{maile_when_2022} proposes two
such triggers. The activation trigger $\phi_a^{\mathrm{ED}}$~(Eq.~\ref{eq:north_main}) monitors the activation effective rank of a layer's
activations $\Postact[l]$ relative to its width $C_l$:
\begin{equation}\label{eq:north_main}
\phi_a^{\mathrm{ED}}(f,l)
:= \frac{1}{C_l}
\left|\left\{\sigma\in\operatorname{SVD}(\Postact[l]/\sqrt n):
\sigma>\varepsilon\right\}\right|.
\end{equation}
While this ratio remains below its value at
initialization, the layer is deemed to retain unused directions in which new
features can still be built with the current architecture. However, since there is no direct dependence on the loss, such activation-based criteria can only measure whether existing capacity is saturated, not whether further capacity would be useful.

Predicting the benefit of capacity that does not yet exist is the central idea of several growing methods. Splitting methods (S2D, S2D-E, S3D) \citep{liu_splitting_2019,wang_energy-aware_2020,wu_steepest_2021} evaluate after convergence whether the network is stuck at a stationary point of the loss landscape that additional capacity would allow it to escape. More explicitly, the \emph{expressivity bottleneck} of \emph{TINY}
\citep{verbockhaven_growing_2024} and the natural expansion increment
$\Delta\eta_{\mathrm{SENN}}$ of \emph{SENN} \citep{mitchell_self-expanding_2024} characterize
the part of the desired functional update that the current layer cannot
express. Note that SENN's $\Delta\eta_{\mathrm{SENN}}$ is compared with an absolute threshold and a relative threshold $\Delta\eta_{\mathrm{SENN}}/\eta_0 > \tau$, where $\eta_0$ is the score of the current network. The use of a relative threshold can be linked to theoretical results \citep{csillagFunctionalGradientDescent2026} that provide convergence guarantees for a network for which the ratio of the bottleneck to the total gradient is small enough.

\subsection{When to stop}

Most growing methods follow a predetermined schedule, providing no notion of when to stop, and only a small number have an explicit (or implicit) notion of stopping, generally based on their \emph{when to grow} criteria. AutoGrow~\citep{wen_autogrow_2020} permanently freezes a sub-network when its validation-accuracy gain over a stopping window, $\Delta\mathrm{Acc}_{\mathrm{val}}$, falls below a threshold. Splitting (S2D, S2D-E, S3D)~\citep{liu_splitting_2019,wang_energy-aware_2020,wu_steepest_2021} methods also naturally terminate when a local minimum of the loss is reached, while grow-prune methods such as SENN~\citep{mitchell_self-expanding_2024} achieve a steady-state behavior by iterative growing and pruning.


\section{Beyond width and depth growth}\label{sec:beyond_width_depth}

This section considers extensions beyond the width and depth growing setup of the previous section: We discuss the theoretical framework of boosting, which shares many similarities with growth, and also discuss non-trivial network topologies such as sparse networks or growing DAGs. We postpone the application to transformer growth until Sec.~\ref{sec:transformers}, where we provide a self-contained treatment, including empirical results.


\subsection{Gradient Boosting as Growth}\label{sec:boosting}

A precursor to modern Neural Network growth, gradient boosting provides a general framework for constructing additive models by iteratively adding weak learners to minimize a given loss function \citep{friedman_greedy_2001}. The weak learners are chosen from a predefined set of functions $\mathcal{H}$. From an optimization perspective, each boosting iteration can be interpreted as a greedy descent step in the functional space spanned by $\mathcal{H}$. The output function $F$ is built as a weighted sum of weak learners:
$$F(x) = \sum_{m=1}^M \omega_m \gamma_m(x)$$
where $\gamma_m$ are the weak learners and $\omega_m$ their corresponding weights. The goal is to minimize a loss function $\ell{}(y, F(x))$ over the training data. At iteration $m$, the weak learner $\gamma_m$ and its weight $\omega_m$ are chosen to minimize the loss:
$$(\gamma_m, \omega_m) = \argmin_{\gamma \in \mathcal{H}, \omega\in\mathbb{R}}  \EE[\ell\left(y, F_{m-1}(x) + \omega \gamma(x)\right)]$$
where $F_{m-1}$ is the current model. 
Boosting algorithms typically select the weak learner that is maximally aligned with the negative gradient of the loss with respect to the current model's output.

In the context of neural network layer width growth, adding a neuron at layer $l-1$ induces a functional perturbation $\deltas$ of the pre-activation at layer $l$, which can be interpreted as a boosting step, though over a continuous set of possible weak learners.
Unlike classical boosting methods, except when growth occurs in the last hidden layer, the functional step is added to an intermediate representation.  
Finding the optimal weights for the newly added neurons to match the desired functional variation of the pre-activation at layer $l$ leads to a non-convex optimization problem that is generally NP-hard  \citep{manurangsi2018computationalcomplexitytrainingrelus,bach_breaking_2017} and thus difficult to scale~\citep{liu_splitting_2019}.

\textbf{AdaNet.} Boosting directly inspires AdaNet~\citep{cortesAdaNetAdaptiveStructural2017}, which uses neural networks as weak learners. AdaNet selects candidates using the risk--complexity objective
\begin{equation}\label{eq:adanet_main}
F(\mathbf{w}) = \widehat R(\mathbf{w}) + \sum_j \Gamma_j |w_j|,
\end{equation}
where \(\widehat R\) is the empirical surrogate risk and \(\Gamma_j\) penalizes the complexity of weak learner \(j\). At each boosting step, new neural networks are trained while previous ones are frozen; the candidate giving the largest reduction in \(F(\mathbf w)\) is added to the ensemble, and the process continues.

\subsection{Sparse growth}\label{sec:sparse_prune_methods}

Sparse masks are frequently employed as a way to select important neurons and prune the rest. However, as discussed in~\citet{dai_incremental_learning_2022}, in order for a model to withstand multiple reduction and growth steps, the performance hit it takes during pruning should be completely recoverable, if not surpassable, by the growth process. The objective of sparse growth is not to reduce the computational cost of inference but rather to increase the network's capabilities while avoiding over-parameterization.

\textbf{Grow-and-Prune.} To perform incremental learning, \citeauthor{dai_incremental_learning_2022} use a gradient-based growth where the gradient of all masked connections is averaged over an epoch, and if it surpasses a specific percentile, they are reactivated. Similarly, they prune or deactivate connections when their weight magnitude is below a specific percentile. This two-step process aims to support long-term learning and outperforms simply training from scratch when new data arrive in both error rate and model size.

\textbf{MorphNet}~\citep{gordon_morphnet_2018} uses a sparsity regularizer to penalize over-parameterization while training and then uniformly expands all layers by scaling their width up to a budget. The sparse training maintains good performance and even improves over simple uniform growth under the same FLOPs, showing the benefit of reduction.

In a similar sparse growth manner \citet{yuan_growing_2021} use masking to start from a very sparse seed architecture and utilize budget-driven sparsity regularization to reduce sparsity, thus growing the network progressively. The method achieves higher accuracy than AutoGrow with smaller models and sparse channels.

\textbf{CompNet}~\citep{lu_CompNet_2018} separately trains and imposes an independently interpretable lasso regularization on the new neurons while optimizing for function-preservation. This sparsity optimization can either be applied to the input or output neurons of the new layer that is inserted in the network.

\subsection{DAG growth}\label{sec:dag_growth}

Most generally, one can consider growing an arbitrary Directed Acyclic Graph (DAG). The structure of a neural network can be represented by a DAG in various ways, depending on the degree of granularity required. In classic works such as NEAT~\citep{stanley_neat_2002}, each vertex represents a single neuron, and the edges represent weights, which are then grown using evolutionary methods. Alternatively, each vertex (or node) can represent an entire layer, as is commonly used as a search space for NAS~\citep{nasbench101}.

The flexibility of a DAG representation does not come without trade-offs, as the number of potential growth operations increases quadratically with the vertex and edge count. To improve on the brute-force strategy of growing every possible node, \citet{douka_growth_2025} propose selecting the node $v \in V$ which maximizes the expressivity bottleneck $v^*=\argmax_{v \in V} \frob{\Bottleneck_v}$, and then only considering growth operations which influence this node $v$, either by expanding that node or adding edges connected to it.


\section{Non-stationary data distributions}\label{sec:non_stationary}

A natural application of neuron addition is to non-stationary data distributions, such as Continual Learning. Instead of a single task, represented by its dataset $\mathcal{D}$, we have a sequence of tasks $\mathcal{D}_1, \dots, \mathcal{D}_T$, and want a model that performs well on all tasks $\{\mathcal{D}_1, \dots, \mathcal{D}_T\}$. In addition to the objective of maximising performance on the current task $\mathcal{D}_i$, the goal is also to prevent relapse on previous tasks, known as \emph{catastrophic forgetting}~\citep{lange_continual_2021}. The trade-off between performing well on the new task while retaining performance on the old tasks is known as the \emph{plasticity-stability dilemma}. Parameter isolation methods, which dedicate specific parameters to individual tasks, are a particularly strong family of Continual Learning algorithms, and Neural Network growth is naturally adapted to this, as in Progressive Neural Networks~\citep{rusu_progressive_2016}, which add new layers for each task.

\textbf{Continual Learning.} By aligning the architectures and weights $(A_t, \theta_t)$ with the current task $\D_t$, growing allows one to reuse existing weights and only add parameters for the next task $\D_{t+1}$ when necessary. Some methods enforce sparsity in the weights ~\citep{yoon_lifelong_2018,yang_grown_2021} or temporary pruning ~\citep{hung_CPG_2019,wu_firefly_2020} --- disabling the weights to be able to adapt them when training the next task without introducing semantic drift. This results in supernets where each task uses a subset of the weights. Based on ablation studies~\citep{yoon_lifelong_2018}, growth avoids catastrophic forgetting and adds flexibility to the parameters while maintaining smaller models. The \textbf{Learn-to-Grow} framework~\citep{li_learn_2019} uses NAS to balance performance and accuracy by reusing existing weights as much as possible and replacing them when necessary.

\textbf{Reinforcement Learning.} Loss of plasticity~\citep{dohare_loss_2024} is also a significant problem in Reinforcement Learning: the value function and/or policy are represented by Neural Networks and quickly overfit to initial observations, failing to adapt to new data as training progresses. Proposed solutions include periodically resetting some or all of the network~\citep{nikishin_primacy_2022,sokar_dormant_2023}, or using various forms of regularization~\citep{bhatt_crossq_2024,galliciSimplifyingDeepTemporal2025}. 
Recently, growth methods have been used to maintain plasticity with a fixed growth schedule~\citep{kangForgetandGrowStrategyDeep2025} or with sparse grow-prune methods~\citep{liu_neuroplastic_2025}.


\section{Empirical results}\label{sec:results}

\makeatletter
\begin{table}
\centering
\caption{Aggregated self-reported results for width-growing methods. $\Delta\textrm{scratch acc.}$ indicates the change in accuracy relative to the final architecture trained from scratch. $\Delta\textrm{baseline acc.}$ indicates the change in test (or validation) accuracy relative to a fixed baseline. $\Delta\textrm{params}$, $\Delta\textrm{final compute}$, $\Delta\textrm{total compute}$, and $\Delta\textrm{runtime}$ indicate changes in final parameter count, final compute, total compute (MACs or FLOPs), and total training time, respectively, expressed as a percentage relative to the baseline.}\label{tab:results_width_growing}
\fontsize{5pt}{6pt}\selectfont
\setlength{\tabcolsep}{1pt}
\noindent
\begin{tabularx}{\textwidth}{>{\centering\arraybackslash}p{0.05\linewidth}p{0.15\linewidth}p{0.10\linewidth}X>{\centering\arraybackslash}p{0.09\linewidth}>{\centering\arraybackslash}p{0.09\linewidth}>{\centering\arraybackslash}p{0.09\linewidth}>{\centering\arraybackslash}p{0.08\linewidth}>{\centering\arraybackslash}p{0.08\linewidth}>{\centering\arraybackslash}p{0.09\linewidth}}
\toprule
\multicolumn{1}{>{\centering\arraybackslash}m{0.05\linewidth}}{} &
\multicolumn{1}{m{0.15\linewidth}}{Method} &
\multicolumn{1}{m{0.10\linewidth}}{Dataset} &
\multicolumn{1}{m{\TX@col@width}}{Baseline\newline Architecture} &
\multicolumn{1}{>{\centering\arraybackslash}m{0.09\linewidth}}{$\Delta$ scratch acc. (pp)} &
\multicolumn{1}{>{\centering\arraybackslash}m{0.09\linewidth}}{$\Delta$ baseline acc. (pp)} &
\multicolumn{1}{>{\centering\arraybackslash}m{0.09\linewidth}}{$\Delta$ params (\%)} &
\multicolumn{1}{>{\centering\arraybackslash}m{0.08\linewidth}}{$\Delta$ final compute (\%)} &
\multicolumn{1}{>{\centering\arraybackslash}m{0.08\linewidth}}{$\Delta$ total compute (\%)} &
\multicolumn{1}{>{\centering\arraybackslash}m{0.09\linewidth}}{$\Delta$ runtime (\%)} \\
\midrule
\multirow[c]{18}{=}{\centering\rotatebox[origin=c]{90}{\scriptsize Fixed final arch.}}
 & Net2Net & ImageNet & Inception-BN      & -0.1\sdag\sddag & --              & 0.0          & 0.0 & -- & -- \\
\cmidrule{2-10}
 & GradMax & CIFAR-10  & VGG-11       & -2.2 & -- & 0.0 & 0.0 & -- & -- \\
 & GradMax & CIFAR-10  & WRN-28-1     & -1.8 & -- & 0.0 & 0.0 & -- & -- \\
 & GradMax & CIFAR-100 & WRN-28-1     & -2.5 & -- & 0.0 & 0.0 & -- & -- \\
 & GradMax & ImageNet  & MobileNet-V1 & -2.2 & -- & 0.0 & 0.0 & -- & -- \\
\cmidrule{2-10}
 & TINY & CIFAR-100 & ResNet-18 & -2.5 & -- & 0.0 & 0.0 & 0.0  & -- \\
\cmidrule{2-10}
 & Variance Transfer & CIFAR-10  & VGG-11      & +0.2 & -- & 0.0 & 0.0 & -47.1 & -- \\
 & Variance Transfer & CIFAR-10  & ResNet-20   & -0.1 & -- & 0.0 & 0.0 & -45.1 & -- \\
 & Variance Transfer & CIFAR-10  & MobileNetV1 & -0.3 & -- & 0.0 & 0.0 & -46.2 & -- \\
 & Variance Transfer & CIFAR-100 & VGG-19      & +0.7 & -- & 0.0 & 0.0 & -47.9 & -- \\
 & Variance Transfer & CIFAR-100 & ResNet-18   & -0.2 & -- & 0.0 & 0.0 & -47.4 & -14.3 \\
 & Variance Transfer & CIFAR-100 & MobileNetV1 & -0.6 & -- & 0.0 & 0.0 & -47.1 & -- \\
 & Variance Transfer & ImageNet  & ResNet-50   & -0.8 & -- & 0.0 & 0.0 & -39.9 & -- \\
 & Variance Transfer & ImageNet  & MobileNetV1 & -0.9 & -- & 0.0 & 0.0 & -36.3 & -- \\
\midrule
\multirow[c]{11}{=}{\centering\rotatebox[origin=c]{90}{\scriptsize Adaptive}}
 & Splitting S2D & CIFAR-10 & VGG-19      & -- & -0.1\sdag\sddag & -85.0\sdag & -- & -- & -- \\
 & Splitting S2D & CIFAR-10 & MobileNetV1 & -- & +1.2\sdag\sddag & -88.0\sdag & -- & -- & -- \\
\cmidrule{2-10}
 & Splitting S2D-E & CIFAR-100 & MobileNetV1   & -- & +1.7\sdag\sddag & -- & -34.0\sdag & -- & -- \\
 & Splitting S2D-E & ImageNet  & MobileNetV1   & -- & +1.0\sddag & -- & -1.4  & -- & -- \\
 & Splitting S2D-E & ImageNet  & MobileNetV2   & -- & +0.8\sddag & -- & -0.7  & -- & -- \\
\cmidrule{2-10}
 & Splitting S3D   & CIFAR-100 & MobileNetV1   & -- & +1.2\ddag & -90.9 & -21.7 & -- & -- \\
 & Splitting S3D   & ImageNet  & MobileNetV1   & -- & +1.2\ddag & --    & -1.9  & -- & -- \\
 & NORTH           & CIFAR-10  & VGG-11        & -- & +0.4\sdag & -65.9\sdag & -- & -- & -- \\
 & NORTH           & CIFAR-10  & WideResNet-28 & -- & -0.3\sdag & -16.8\sdag & -- & -- & -- \\
 & NORTH           & CIFAR-100 & VGG-11        & -- & +0.8\sdag & -68.0\sdag & -- & -- & -- \\
\bottomrule
\end{tabularx}
\begin{tablenotes}[flushleft]
\footnotesize
\item[\sdag] Results derived from plots rather than tabulated values.
\item[\sddag] Grown networks are optimized for significantly more gradient steps than the baseline.
\end{tablenotes}
\end{table}

\begin{table}
\centering
\caption{Aggregated self-reported results for depth-growing methods. See the caption of Table~\ref{tab:results_width_growing} for definitions of each metric. }\label{tab:results_depth_growing}
\fontsize{5pt}{6pt}\selectfont
\setlength{\tabcolsep}{1pt}
\noindent
\begin{tabularx}{\textwidth}{>{\centering\arraybackslash}p{0.05\linewidth}p{0.11\linewidth}p{0.14\linewidth}X>{\centering\arraybackslash}p{0.10\linewidth}>{\centering\arraybackslash}p{0.10\linewidth}>{\centering\arraybackslash}p{0.07\linewidth}>{\centering\arraybackslash}p{0.07\linewidth}>{\centering\arraybackslash}p{0.07\linewidth}>{\centering\arraybackslash}p{0.07\linewidth}}
\toprule
\multicolumn{1}{>{\centering\arraybackslash}m{0.05\linewidth}}{} &
\multicolumn{1}{m{0.11\linewidth}}{Method} &
\multicolumn{1}{m{0.14\linewidth}}{Dataset} &
\multicolumn{1}{m{\TX@col@width}}{Baseline\newline Architecture} &
\multicolumn{1}{>{\centering\arraybackslash}m{0.10\linewidth}}{$\Delta$ scratch acc. (pp)} &
\multicolumn{1}{>{\centering\arraybackslash}m{0.10\linewidth}}{$\Delta$ baseline acc. (pp)} &
\multicolumn{1}{>{\centering\arraybackslash}m{0.07\linewidth}}{$\Delta$ params (\%)} &
\multicolumn{1}{>{\centering\arraybackslash}m{0.07\linewidth}}{$\Delta$ final compute (\%)} &
\multicolumn{1}{>{\centering\arraybackslash}m{0.07\linewidth}}{$\Delta$ total compute (\%)} &
\multicolumn{1}{>{\centering\arraybackslash}m{0.07\linewidth}}{$\Delta$ runtime (\%)} \\
\midrule
\multirow[c]{16}{=}{\centering\rotatebox[origin=c]{90}{\scriptsize Fixed final arch.}}
 & Net2Net & ImageNet & Inception-BN & 0.0\sdag\sddag & -- & 0.0\sddag & -- & -- & -- \\
\cmidrule{2-10}
 & LipGrow & CIFAR-10     & ResNet-50 & -0.6 & -- & 0.0 & -- & -56.6${}^*$ & -- \\
 & LipGrow & CIFAR-10     & ResNet-74 & -0.3 & -- & 0.0 & -- & -53.0${}^*$ & -44.8\sdag \\
 & LipGrow & CIFAR-100    & ResNet-50 & +0.3 & -- & 0.0 & -- & -46.1${}^*$ & -- \\
 & LipGrow & CIFAR-100    & ResNet-74 & -0.4 & -- & 0.0 & -- & -53.0${}^*$ & -- \\
 & LipGrow & TinyImageNet & ResNet-66 & +0.7 & -- & 0.0 & -- & -47.8${}^*$ & -- \\
\cmidrule{2-10}
 & FRAGrow & CIFAR-10 & VGG           & 0.0  & -- & 0.0 & -- & --    & -18.4 \\
 & FRAGrow & CIFAR-10 & ResNet        & +0.3 & -- & 0.0 & -- & --    & -17.6 \\
 & FRAGrow & CIFAR-10 & MobileNetV2   & -0.4 & -- & 0.0 & -- & --    & -20.6 \\
 & FRAGrow & CIFAR-100 & VGG          & +0.4 & -- & 0.0 & -- & --    & -16.0 \\
 & FRAGrow & CIFAR-100 & ResNet       & +0.4 & -- & 0.0 & -- & --    & -20.1 \\
 & FRAGrow & CIFAR-100 & MobileNetV2  & 0.0  & -- & 0.0 & -- & --    & -26.2 \\
 & FRAGrow & ImageNet & VGG           & -0.2 & -- & 0.0 & -- & --    & -11.2 \\
 & FRAGrow & ImageNet & ResNet        & -0.2 & -- & 0.0 & -- & --    & -10.0 \\
 & FRAGrow & ImageNet & MobileNetV2   & -0.5 & -- & 0.0 & -- & --    & -6.3  \\
\midrule
\multirow[c]{8}{=}{\centering\rotatebox[origin=c]{90}{\scriptsize Adaptive}}
 & AutoGrow & CIFAR-10  & ResNet-56 (3 stage)          & 0.0   & +0.4 & +374.9      & -- & --          & -- \\
 & AutoGrow & CIFAR-10  & ResNet (4 stage)             & -0.1  & --   & --          & -- & --          & -- \\
 & AutoGrow & CIFAR-10  & PlainNet (3 stage)           & +6.5  & --   & --          & -- & --          & -- \\
 & AutoGrow & CIFAR-10  & PlainNet (4 stage)           & +5.7  & --   & --          & -- & --          & -- \\
 & AutoGrow & CIFAR-100 & ResNet (3 stage)             & -1.0  & --   & --          & -- & --          & -- \\
 & AutoGrow & CIFAR-100 & ResNet (4 stage)             & +1.2  & --   & --          & -- & --          & -- \\
 & AutoGrow & CIFAR-100 & PlainNet (3 stage)           & +31.5 & --   & --          & -- & --          & -- \\
 & AutoGrow & CIFAR-100 & PlainNet (4 stage)           & +29.3 & --   & --          & -- & --          & -- \\
 & AutoGrow & ImageNet  & ResNet-34 (4 stage)          & +0.4  & +3.0\sdag & +214.2\sdag & +198.1\sdag & -- & -- \\
 & AutoGrow & ImageNet  & PlainNet (4 stage)           & +0.7  & --   & --          & -- & --          & -- \\
\bottomrule
\end{tabularx}
\begin{tablenotes}[flushleft]
\footnotesize
\item[\sdag] Results derived from plots rather than tabulated values.
\item[\sddag] Grown networks are optimized for significantly more gradient steps than the baseline.
\item[${}^*$] LipGrow's $\Delta$total compute is reported using the non-standard PPE (parameters-per-epoch) metric, but for the specific doubling architecture of LipGrow this is almost identical to the change in total MACs.
\end{tablenotes}
\end{table}
\makeatother

In addition to applications to specific contexts such as non-stationary data distributions, the most commonly cited motivations in the supervised learning setting are:
\begin{itemize}
    \item \textbf{Optimization benefits}, e.g., due to regularisation: growth may improve final performance relative to training the final architecture from scratch.
    \item \textbf{Frugality} in terms of reduced runtime, computational cost, or final parameter count.
    \item \textbf{Architecture adaptation}, either to the task or offering an adjustable trade-off between performance and some notion of e.g. complexity or computational cost.
\end{itemize}
Here we re-examine the key growing methods, and ask whether they do in fact satisfy these motivations.
We aggregate the self-reported results for the largest classes of growing algorithms, width growing (Table~\ref{tab:results_width_growing}) and depth growing (Table~\ref{tab:results_depth_growing}). We include all surveyed methods that grow exclusively width or depth on image classification (CIFAR-10, CIFAR-100, TinyImageNet, ImageNet), and that provide both numerical baseline references and numerical entries compatible with our set of reported metrics.

To reduce the impact of differences in hyperparameters and experimental setups, we report only the relative changes in the following metrics: $\Delta \textrm{scratch acc.}$ indicates the percentage point change in accuracy relative to retraining the final architecture from scratch. $\Delta \textrm{baseline acc.}$ indicates the change in accuracy relative to the baseline architecture indicated in the Baseline Architecture column. This is only relevant for the adaptive architecture methods, which may compare to a reference architecture different from the final one. $\Delta \textrm{params}$, $\Delta\textrm{final compute}$ and $\Delta\textrm{total compute}$ indicate changes in final parameter count, final compute and total compute (MACs or FLOPs), respectively, while $\Delta\textrm{runtime}$ indicates total training runtime. All of these figures are expressed as a percentage relative to the baseline.

\subsection{Optimization benefits}
By comparing the performance of the final grown model to that of the same architecture trained from scratch (without growth), $\Delta\textrm{scratch acc.}$ provides a relatively clean ablation for the optimization or regularisation effects of growth. Width-growing to a fixed final architecture rarely yields a positive $\Delta\textrm{scratch acc.}$ improvement. Net2Net, the oldest method, trains for a total of twice the epochs of the baseline, which likely contributes to its relatively strong performance compared to newer methods, which all keep the number of epochs fixed. Later methods generally improve on Net2Net in a matched-epoch setting. In general, width-growing is only able to improve from-scratch performance on older or harder-to-optimize architectures such as Inception and VGG.

This is also visible in the depth-growing results~\autoref{tab:results_depth_growing}, where the most significant $\Delta\textrm{scratch acc.}$ results are for the PlainNet architectures (ResNet architectures without the residual connections).
In general, depth-growing can match and occasionally slightly improve performance, even with newer ResNet architectures. Such gains have been attributed to the regularising effect of growth~\cite{wu_when_2024}. This is consistent with the depth-growing discussion of Sec.~\ref{sec:where_when}: growth timing can act as a regularizer to control over/underfitting.

\subsection{Frugality}
To gauge frugality, we compare the change in final parameters ($\Delta\textrm{params}$), final and total compute ($\Delta\textrm{final compute}$ and $\Delta\textrm{total compute}$), and total training runtime ($\Delta\textrm{runtime}$). 

For width-growing convolutional networks, FLOPs and MACs scale with both input and output channels, so reduced channel counts can produce large $\Delta\textrm{final compute}$ reductions. However, this does not straightforwardly translate into reductions in runtime (or energy), as parameters such as the batch size are tuned for the large models and need to be adapted to the smaller one. As a case study, Variance Transfer~\citep{yuan_accelerated_2023} observes a mere 14.3\% runtime reduction when growing a ResNet-18 on CIFAR-100, despite reducing total FLOPs by 47.4\% (see Table~\ref{tab:results_width_growing}). As a remedy, they propose adapting the batch size to maximize GPU memory usage and achieve a 40\% runtime reduction for a negligible change in final accuracy. However, large-batch training often results in poor generalization performance~\citep{keskar_large-batch_2017} and therefore may not be generally applicable.

In contrast, since compute and runtime scale approximately linearly with the number of layers, depth growth translates more directly into runtime savings (up to 45\% in Table~\ref{tab:results_depth_growing}).

Although many methods report accuracy metrics, they often do not comprehensively report training times. This is the case of e.g. GradMax~\citep{evci_gradmax_2022}, despite explicitly aiming to reduce training time. Such missing timings are indicated by a dash.

\subsection{Architecture adaptation}

The empirical evidence for growth as a form of architecture adaptation is much less complete: for width-growth, the splitting family of methods~\citep{liu_splitting_2019,wang_energy-aware_2020,wu_steepest_2021} is able to find significantly smaller architectures that outperform the standard baselines, as measured by $\Delta \textrm{baseline acc.}$, suggesting that the channel widths in the standard architectures benefit from task-specific tuning. However, as they do not also compare training from scratch ($\Delta \textrm{scratch acc.}$), it is difficult to disentangle from the optimization effects of growth, especially as they retrain for the total number of epochs after each split, leading to a much larger total training budget than the baselines.

For depth-growth, few comparisons to standard architectures are available. AutoGrow~\citep{wen_autogrow_2020} achieves improved performance with a similar or reduced training compute, but the final architecture has very large parameter counts. By comparing with the depth-growing methods with a fixed final architecture, this suggests that while it benefits from the computational savings of depth growth, the final architecture found does not appear to improve in terms of performance-complexity over the standard hand-designed ResNet family. However, a more thorough experimental evaluation of adaptive depth growing methods would be needed to conclude this decisively.


\subsection{Transformer results}\label{sec:results_transformers}

Transformers provide large-scale empirical validation of growing methods.
In general, the motivation is frugality: reduced training cost with a fixed target architecture, sometimes termed \emph{progressive learning} in this context~\citep{li_autoprog_2022,yao_msg_2024}.
These studies test whether the empirical conclusions of the previous section generalize to the billion-parameter scale.

\textbf{What works.} The large-scale study~\citep{du_stacking_2024} demonstrates that depthwise stacking reduces the number of training tokens needed to reach the same loss by 35.3\% when training a 7B parameter model. Compound growth~\citep{gu_bert_2021}, expanding depth, width, and sequence length simultaneously, reduces BERT training time by 42.4--45.1\%. At smaller scales, MSG~\citep{yao_msg_2024} reduces BERT-Large training time by 54.5\%, and AutoProg~\citep{li_autoprog_2022} reduces Vision Transformer training time by up to 46.0\%. These results establish growing methods as a practical, frugal alternative to training from scratch in the evaluated settings.

Similar to the width growing results of the previous section, \citet{du_stacking_2024} report modest speedups from width growth alone. However, recent results from SPARKLING show that width growth can be effective if performed sufficiently late during training, indicating the importance of \emph{when to grow} in the transformer setting~\citep{yu_sparkling_2026}.

\textbf{Function preservation.} MSG~\citep{yao_msg_2024} achieves strong results using strict function preservation, while AutoProg's ReZero-based identity-insertion ablation~\citep{li_autoprog_2022} reduces Vision Transformer performance by 3.21 percentage points relative to the fixed-model baseline. Most strikingly, simple stacking~\citep{du_stacking_2024}--which is not function-preserving--achieves the best results among the operators in the 1.1B comparison and remains effective at 7B.

\textbf{Optimization benefits.} Growth can also have optimization benefits beyond a reduction in training cost.
Stacking studies report improved task performance, particularly on reasoning tasks, even at matched perplexity~\citep{saunshi_inductive_2024,kapl2025depthgrownmodelsovercomecurse}.
One proposed explanation concerns the ``curse of depth'', whereby many layers in the latter half of a transformer contribute little to its output~\citep{sun_curse_of_depth_2025}, and it has been suggested that depth growth mitigates this behavior~\citep{kapl2025depthgrownmodelsovercomecurse}.

These experiments validate growing methods as a frugal training strategy at scale, but expose limits in current understanding. A significant scale gap remains: at the time of writing, the majority of the experiments are at most 8B parameters, while production language models exceed 70B.


\section{Open Questions and Future Directions}\label{sec:conclusion}

\input{conclusion.tex}

\backmatter

\section*{Acknowledgements}

This work was supported by grants ANR-22-CE33-0015-01 operated by LISN to S.~Chevallier and by ANR-20-CE23-0025 operated by Inria to G.~Charpiat. 
The EU also funded this work under GA no.101135782 (MANOLO project).

\section*{Contributions}

\begin{itemize}
    \item Alex Davey -- Overall editing, introduction, conclusion and empirical results
    \item Théo Rudkiewicz -- Overall editing, Unified framework, function geometry and when to grow section
    \item Stéphane Rivaud -- Transformer and depth growth section
    \item Stella Douka -- Continual Learning and Sparse Grow-Prune
    \item Félix Houdouin -- Gradient Boosting
    \item Julien Mille -- Proofreading
    \item Alena Shilova -- Proofreading
    \item Guillaume Charpiat -- Proofreading
    \item Sylvain Chevallier -- Proofreading and original idea conception
\end{itemize}


\newpage
\begin{appendices}
\section{Proofs}\label{app:proofs}

This appendix presents the proofs for the unified framework \autoref{thm:unified_obj_main_text}. We show that neuron addition methods
such as TINY~\citep{verbockhaven_growing_2024}, GradMax~\citep{evci_gradmax_2022}, SENN~\citep{mitchell_self-expanding_2024}, and NeST~\citep{daiNeSTNeuralNetwork2019} can all be understood as special
cases of a common optimization objective, differing only in their choice
of metric and constraints.

\subsection{Preliminaries}

\begin{lemma}\label{lemma:argument_scaling}

For:

\begin{itemize}
\item \(E\) a pre-Hilbert space,
\item \(F\) a non empty set stable by scaling \emph{i.e.}
  (\(\forall x \in F,\forall\lambda \in \mathbb{R},\lambda x \in F\))
\item \(f:F \rightarrow E\) respect some scaling homogeneity
  (\(\exists p \in \mathbb{N}^{\ast},\forall x \in F,\forall\lambda \in \mathbb{R},f(\lambda x) = \lambda^{p}f(x)\))
\item \(y \in E\).
\end{itemize}
If \(\left( \exists x \in F,{\langle f(x),y\rangle} > 0 \right)\) we
have that:
\[\argmax_{d \in F~|~\left\| {f(d)} \right\| \leq 1}{\langle y,f(d)\rangle} = \argmax_{d \in F~|~\left\| {f(d)} \right\| = 1}{\langle y,f(d)\rangle}\]
and:
\[\argmin_{x \in F}\left\| {y - f(x)} \right\| = \argmin_{x \in F~|~\left\| {f(x)} \right\| \neq 0}\left\| {y - f(x)} \right\|\]

\end{lemma}

\begin{proof}

Let
\(I \coloneqq \argmax_{d \in F|\left\| {f(d)} \right\| \leq 1}{\langle y,f(d)\rangle}\)
and
\(R \coloneqq \argmax_{d \in F|\left\| {f(d)} \right\| = 1}{\langle y,f(d)\rangle}\).

\begin{itemize}
\item
  Let us show that \(I \subset R\). Let \(d^{\ast} \in I\); we will show that
  \(\left\| {f\left( d^{\ast} \right)} \right\| = 1\).

  Let \(x \in F\) with \({\langle f(x),y\rangle} > 0\) (exists by
  hypothesis). For \(\varepsilon > 0\) small enough,
  \(\left\| {f(\varepsilon x)} \right\| = \varepsilon^{p}\left\| {f(x)} \right\| \leq 1\),
  so \(\varepsilon x\) satisfies the norm constraint in the definition of \(I\), and
  \({\langle y,f(\varepsilon x)\rangle} = \varepsilon^{p}{\langle y,f(x)\rangle} > 0\).
  By optimality of \(d^{\ast}\) in \(I\),
  \({\langle y,f\left( d^{\ast} \right)\rangle} \geq {\langle y,f(\varepsilon x)\rangle} > 0\),
  hence \(f\left( d^{\ast} \right) \neq 0\). Let
  \(\lambda \coloneqq \frac{1}{\left\| {f\left( d^{\ast} \right)} \right\|^{\frac{1}{p}}} > 0\).
  As \(|\lambda| \geq 1\):
  \[{\langle y,f\left( d^{\ast} \right)\rangle} \leq \lambda^{p}{\langle y,f\left( d^{\ast} \right)\rangle} = {\langle y,f\left( \lambda d^{\ast} \right)\rangle}\]
  by homogeneity of \(f\). By maximality of \(d^{\ast}\) in \(I\), as
  \(\left\| {f\left( \lambda d^{\ast} \right)} \right\| = 1\), this is
  an equality hence \(\lambda = 1\) and
  \(\left\| {f\left( d^{\ast} \right)} \right\| = 1\). Therefore
  \(d^{\ast} \in R\) and \(I \subset R\).
\item
  Let us show that \(R \subset I\). Let \(d^{\ast} \in R\); since
  \(I \subset R\), \(d^{\ast}\) maximizes the score over
  \(\left\{ d \in F~|~\left\| {f(d)} \right\| \leq 1 \right\}\), and as
  \(\left\| {f\left( d^{\ast} \right)} \right\| \leq 1\) we have that
  \(d^{\ast} \in I\). We conclude that \(R \subset I\).
\end{itemize}

For the second equality, with the witness \(x\) above,
\(\left\| {y - f(\varepsilon x)} \right\|^{2} = \left\| y \right\|^{2} - 2\varepsilon^{p}{\langle y,f(x)\rangle} + \varepsilon^{2p}\left\| {f(x)} \right\|^{2} < \left\| y \right\|^{2}\)
for \(\varepsilon > 0\) small, while \(f(x') = 0\) gives
\(\left\| {y - f(x')} \right\|^{2} = \left\| y \right\|^{2}\); so no
minimizer of \(\left\| {y - f(x)} \right\|\) has \(f(x) = 0\).

\end{proof}

\begin{lemma}\label{lemma:sign_symmetrization}

For:

\begin{itemize}
\item
  \(E\) a pre-Hilbert space,
\item
  \(F\) a non-empty set,
\item
  \(f:F \rightarrow E\),
\item
  \(y \in E\),
\item
  \(n:F \rightarrow F\) a sign-flip operator:
  \(\forall d \in F,f\left( n(d) \right) = - f(d)\).
\end{itemize}

Then \(\langle f\left( . \right),y\rangle\) takes a strictly positive
value as soon as it takes a non-zero one:
\[\left( \exists x \in F,{\langle f(x),y\rangle} \neq 0 \right) \Longrightarrow \left( \exists x \in F,{\langle f(x),y\rangle} > 0 \right).\]

\end{lemma}

\begin{proof}

Let \(x \in F\) with \({\langle f(x),y\rangle} \neq 0\). If
\({\langle f(x),y\rangle} > 0\) we are done; otherwise
\({\langle f(x),y\rangle} < 0\) and \(n(x) \in F\) satisfies
\({\langle f\left( n(x) \right),y\rangle} = - {\langle f(x),y\rangle} > 0\).

\end{proof}

\begin{theorem}\label{lemma:argmin_norm_to_argmax_scalarp}

For:

\begin{itemize}
\item
  \(E\) a pre-Hilbert space,
\item
  \(F\) a non empty set stable by scaling \emph{i.e.}
  (\(\forall x \in F,\forall\lambda \in \mathbb{R},\lambda x \in F\))
\item
  \(f:F \rightarrow E\) respect some scaling homogeneity
  (\(\exists p \in \mathbb{N}^{\ast},\forall x \in F,\forall\lambda \in \mathbb{R},f(\lambda x) = \lambda^{p}f(x)\))
\item
  \(n:F \rightarrow F\) a sign-flip operator:
  \(\forall d \in F,f\left( n(d) \right) = - f(d)\)
\item
  \(y \in E\) such that
  \(\exists x \in F,{\langle f(x),y\rangle} \neq 0\) (non-degeneracy).
\end{itemize}

We have that:

\begin{equation}\label{eq:fixed_norm_optimization_equivalence}
\argmin_{x \in F}\left\| {y - f(x)} \right\| = \left\{ {\langle y,f\left( d^{\ast} \right)\rangle}^{\frac{1}{p}}d^{\ast}~|~d^{\ast} \in \argmax_{d \in F~|~\left\| {f(d)} \right\| \leq 1}{\langle y,f(d)\rangle} \right\}
\end{equation}

\end{theorem}

\begin{proof}

By \autoref{lemma:sign_symmetrization} the non-degeneracy hypothesis yields
some \(x \in F\) with \({\langle f(x),y\rangle} > 0\), so by
\autoref{lemma:argument_scaling} we can add the constraint
\(\left\| {f(x)} \right\| \neq 0\) in the left-hand side and replace
\(\leq 1\) by \(= 1\) in the right-hand side. Let's proceed by setting
equality. \[\begin{aligned}
\argmin_{x \in F~|~\left\| {f(x)} \right\| \neq 0}\left\| {y - f(x)} \right\| & = \argmin_{x \in F}\left( \left\| y \right\|^{2} - 2{\langle y,f(x)\rangle}+ \left\| {f(x)} \right\|^{2}  \right) \\
 & = \argmin_{x \in F~|~\left\| {f(x)} \right\| \neq 0}\left( - 2{\langle y,f(x)\rangle} + \left\| {f(x)} \right\|^{2} \right) \\
 & = \argmin_{x \in F~|~\left\| {f(x)} \right\| \neq 0}\left( - 2{\langle y,\left\| {f(x)} \right\|\frac{f(x)}{\left\| {f(x)} \right\|}\rangle} + \left\| {f(x)} \right\|^{2} \right) \\
\text{As}\ f\ \text{is}\ \text{homogenous}: & = \argmin_{x \in F~|~\left\| {f(x)} \right\| \neq 0}\left( - 2\left\| {f(x)} \right\|{\langle y,f\left( \frac{x}{\left\| {f(x)} \right\|^{\frac{1}{p}}} \right)\rangle} + \left\| {f(x)} \right\|^{2} \right) \\
 & = \argmin_{x \in F~|~\left\| {f(x)} \right\| \neq 0}\left( - 2t{\langle y,f(d)\rangle} + t^{2} \right)
\end{aligned}\] with
\(d \coloneqq \frac{x}{\left\| {f(x)} \right\|^{\frac{1}{p}}} \in F\)
and \(t \coloneqq \left\| {f(x)} \right\| \geq 0\). The solution over
\(t\) is achieved at \(t = {\langle y,f(d)\rangle}\). Hence we get:
\[\begin{aligned}
\argmin_{x \in F~|~\left\| {f(x)} \right\| \neq 0}\left\| {y - f(x)} \right\| & = \argmin_{x \in F~|~\left\| {f(x)} \right\| \neq 0}\left( - {\langle y,f(d)\rangle}^{2} \right) \\
 & = \argmax_{x \in F~|~\left\| {f(x)} \right\| \neq 0}{\langle y,f(d)\rangle} \\
 & = \left\{ {\langle y,f\left( d^{\ast} \right)\rangle}^{\frac{1}{p}}d^{\ast}~|~d^{\ast} \in \argmax_{d \in F~|~\left\| {f(d)} \right\| = 1}{\langle y,f(d)\rangle} \right\}
\end{aligned}\]

\end{proof}

\begin{corollary}\label{corollary:argmin_norm_to_argmax_scalarp}

For:

\begin{itemize}
\item
  \(E\) a pre-Hilbert space,
\item
  \(F\) a non empty set stable by scaling \emph{i.e.}
  (\(\forall x \in F,\forall\lambda \in \mathbb{R},\lambda x \in F\))
\item
  \(f:F \rightarrow E\) respect some scaling homogeneity
  (\(\exists p \in \mathbb{N}^{\ast},\forall x \in F,\forall\lambda \in \mathbb{R},f(\lambda x) = \lambda^{p}f(x)\))
\item
  \(n:F \rightarrow F\) a sign-flip operator:
  \(\forall d \in F,f\left( n(d) \right) = - f(d)\)
\item
  \(y \in E\) such that
  \(\left( \exists x \in F,{\langle f(x),y\rangle} \neq 0 \right)\) % or
%  \(\ker f = \left\{ 0 \right\}\).
\end{itemize}

We have that:
\[\argmin_{x \in F}\left\| {y - f(x)} \right\| \propto \argmax_{d \in F|\left\| {f(d)} \right\| \leq 1}{\langle y,f(d)\rangle}\]

\end{corollary}

\begin{proof}
\leavevmode\par\vspace{0.5\baselineskip}

If \(\exists x \in F,{\langle f(x),y\rangle} \neq 0\), the non-degeneracy hypothesis of \autoref{lemma:argmin_norm_to_argmax_scalarp} is met (using the same \(n\)), and its right-hand side is exactly the desired rescaled right-hand \(\argmax\).

% \begin{itemize}
%    \item Otherwise \(\forall x \in F,{\langle f(x),y\rangle} = 0\) and \(\ker f = \left\{ 0 \right\}\). Then \(\left\| {y - f(x)} \right\|^{2} = \left\| y \right\|^{2} + \left\| {f(x)} \right\|^{2} \geq \left\| y \right\|^{2}\), with equality iff \(f(x) = 0\), hence \(\argmin_{x \in F}\left\| {y - f(x)} \right\| = \ker f = \left\{ 0 \right\}\). On the other side \(\max\limits_{\left\| {f(d)} \right\| \leq 1}{\langle y,f(d)\rangle} = 0\), so every maximizer \(d^{\ast}\) satisfies \({\langle y,f\left( d^{\ast} \right)\rangle} = 0\) and the rescaled set \(\left\{ {\langle y,f\left( d^{\ast} \right)\rangle}^{\frac{1}{p}}d^{\ast} \right\} = \left\{ 0 \right\}\). Both sides reduce to \(\left\{ 0 \right\}\), so the proportionality holds.
% \end{itemize}

\end{proof}

\begin{lemma}\label{lemma:optimal_projection}

For:

\begin{itemize}
\item
  \(n,p,a,b,p \in \mathbb{N}\)
\item
  \(Y \in (n,p),X \in (n,a),A \in (a,b),B \in (b,p)\)
\item
  \(s(A):A \mapsto \argmin_{B}\left\| {Y - XAB} \right\|\)
\item
  \(\overline{s}:A \mapsto \overline{s}(A) \in s(A)\)
\item
  \(\hat{Y} \coloneqq XA\overline{s}(A)\)
\end{itemize}

\[\left\| {Y - \hat{Y}} \right\|^{2} = \left\| Y \right\|^{2} - \left\| \hat{Y} \right\|^{2}\]

\end{lemma}

\begin{proof}

By definition of \(\overline{s}(A)\) being optimal, we have that
\(\hat{Y}\) is the orthogonal projection of \(Y\) onto the space spanned
by the columns of \(XA\). Hence \(Y - \hat{Y}\) is orthogonal to
\(\hat{Y}\) and the Pythagorean theorem gives the result.

\end{proof}

\begin{theorem}\label{theorem:separable_optimization}

For:

\begin{itemize}
\item
  \(n,p,a,b,p \in \mathbb{N}\)
\item
  \(Y \in (n,p),X \in (n,a),A \in (a,b),B \in (b,p)\)
\item
  \(s(A):A \mapsto \argmin_{B}\left\| {Y - XAB} \right\|\)
\item
  \(\overline{s}:A \mapsto \overline{s} \in s(A)\)
\end{itemize}

We have that:
\[\argmin_{A,B}\left\| {Y - XAB} \right\| = \bigcup_{A^{\ast}}\left\{ A^{\ast} \right\} \times s\left( A^{\ast} \right)~|~A^{\ast} \in \argmax_{A}\left\| {XA\overline{s}(A)} \right\|\]

\end{theorem}

\begin{proof}

\[
\begin{aligned}
\argmin_{A,B}\left\| {Y - XAB} \right\| & = \argmin_{A}\left\| {Y - XAB} \right\|,B \in s(A)
\end{aligned}
\]
Hence:
\[
\begin{aligned}
    \argmin_{A}\left\| {Y - XAB} \right\| & = \argmin_{A}\left\| {Y - XA\overline{s}(A)} \right\| \\
    \text{By}\ {}: & = \argmin_{A}\left\| Y \right\|^{2} - \left\| \hat{Y} \right\|^{2} \\
     & = \argmax_{A}\left\| {XA\overline{s}(A)} \right\| \\
     & = \bigcup_{A^{\ast}}\left\{ A^{\ast} \right\} \times s(A^{\ast})~|~A^{\ast} \in \argmax_{A}\left\| {XA\overline{s}(A)} \right\|
\end{aligned}
\]

\end{proof}

\begin{theorem}\label{theorem:rimannian_norm}

For:

\begin{itemize}
\item
  \(\left\| . \right\|\) the Frobenius norm
\item
  \(A \in (n,m)\)
\item
  \(X \in (m,n)\) We have that:
  \[\left\| {AX} \right\| = \left\| {\left( A^{\top}A \right)^{\frac{1}{2}}X} \right\|\]
\end{itemize}

\end{theorem}

\begin{proof}

\[\begin{aligned}
\left\| {AX} \right\|^{2} & = {\langle AX,AX\rangle} \\
 & = {\langle A^{\top}AX,X\rangle} \\
 & = {\langle\left( A^{\top}A \right)^{\frac{1}{2}}\left( A^{\top}A \right)^{\frac{1}{2}}X,X\rangle} \\
 & = \left\| {\left( A^{\top}A \right)^{\frac{1}{2}}X} \right\|^{2}
\end{aligned}\]

\end{proof}

\subsubsection{Notations}

\subsubsection{Linear Algebra}

\begin{itemize}
\item
  \(\langle.,.\rangle\) is a scalar product and \(\left\| . \right\|\)
  the associated norm.
\item
  \(\left\| . \right\|_{2}\) is the Frobenius norm.
\item
  \(\times_{k}\) is to indicate multiplication over a common dimension
  \(k\).
\item
  \(.^{+}\) is the Moore-Penrose pseudoinverse.
  \item To make the size of matrices more readable we use $(a, b)$ instead of $\mathbb{R}^{a \times b}$.
\end{itemize}

\subsubsection{Data and Model}

We consider a data distribution \(\D\), a per-data-point loss function
\(\ell:\hat{y} \times y \mapsto \ell(\hat{y},y)\), and a model
\(f:x \times \theta \mapsto f_{\theta}(x) = \hat{y}\) with parameters
\(\theta\) mapping inputs \(x\) to predictions \(\hat{y}\). The overall
loss is:
\[
    \loss(f_{\theta}) = \EE_{(x,y) \sim \D}\left\lbrack \ell(f_{\theta}(x),y) \right\rbrack
\]

\subsubsection{Network Architecture}

We consider a feedforward neural network with layers indexed by \(l\).
For a data point \(x\), the forward pass through consecutive layers is:
\[
    x \overset{a_{l - 2}}{\xrightarrow{\hspace{2em}}} a_{l - 2}(x) \overset{L_{W_{l - 1}}}{\xrightarrow{\hspace{2em}}} s_{l - 1}(x) \overset{\sigma}{\xrightarrow{\hspace{2em}}} a_{l - 1}(x) \overset{L_{W_{l}}}{\xrightarrow{\hspace{2em}}} s_{l}(x) \overset{\hspace{2em}}{\xrightarrow{\hspace{2em}}} f(x)
\]
For readability, we drop layer indices when unambiguous:
\[
    x \overset{a_{- 2}}{\xrightarrow{\hspace{2em}}} a_{- 2}(x) \overset{L_{W_{- 1}}}{\xrightarrow{\hspace{2em}}} s_{- 1}(x) \overset{\sigma}{\xrightarrow{\hspace{2em}}} a_{- 1}(x) \overset{L_{W}}{\xrightarrow{\hspace{2em}}} s(x) \xrightarrow{\hspace{2em}} f(x)
\]

\subsubsection{Dimensions and Batched Notation}

The size of each layer \(l\) is \(C_{l}\) (written as \(C_{- 2}\),
\(C_{- 1}\), \(C\) for layers \(l - 2\), \(l - 1\), \(l\) respectively).
The number of data points is \(n\). For a single sample:
\(a_{- 2}(x) \in \left( C_{- 2},1 \right) = \left( C_{- 2} \right)\). For \(n\) batched samples, we stack row-wise:
\[
    \mathbf{a}_{- 2} = \begin{pmatrix}
    {a_{- 2}(x_{1})}^{\top} \\
    \ldots \\
    {a_{- 2}(x_{n})}^{\top}
    \end{pmatrix} \in \left( n,C_{- 2} \right)
\]
The linear layer acts as: \[\begin{aligned}
\mathrm{L}_{W}(a_{- 2}(x)) & = W_{l - 1} \times_{C_{- 2}}a_{- 2}(x) \in \left( C_{- 1},1 \right) \\
\mathrm{L}_{W}(\mathbf{\mathbf{a}_{- 2}}) & = \mathbf{a}_{- 2} \times_{C_{- 2}}W^{\top} \in \left( n,C_{- 1} \right)
\end{aligned}\]

\subsubsection{Gradient Notation}

The gradient of the loss with respect to pre-activations, stacked
row-wise: \[\Gradpreact = \begin{pmatrix}
{\nabla_{s_{}(x_{1})}\ell(f\left( x_{1} \right))}^{\top} \\
\ldots \\
{\nabla_{s_{}(x_{n})}\ell(f\left( x_{n} \right))}^{\top}
\end{pmatrix} \in (n,C)\]

\subsubsection{Neuron Addition}

We aim to expand the intermediate representation \(s_{- 1}\) and
\(a_{- 1}\) by adding \(\cext\) new neurons. This is done by adding
fan-in weights \(\Fanin{} = \begin{pmatrix}
\fanin{}_{1} \\
\ldots \\
\fanin{}_{k}
\end{pmatrix} \in \left( \cext,C_{- 2} \right)\) and fan-out weights
\(\Fanout{} = \begin{pmatrix}
\fanout{}_{1} & \ldots & \fanout{}_{k}
\end{pmatrix} \in \left( C,\cext \right)\):

\[
a_{- 2}(x) \overset{L_{\Fanin{}}}{\xrightarrow{\hspace{2em}}} \preactext(x) \overset{\sigma}{\xrightarrow{\hspace{2em}}} \postactext(x) \overset{L_{\Fanout{}}}{\xrightarrow{\hspace{2em}}} \delta_s(x)
\]
The change to the pre-activation at layer \(l\) is:
\[s_{} \leftarrow s_{} + \delta_{s_{}}\] with
\(\delta_{s_{}} = \mathrm{L}_{\Fanout{}} \circ \sigma \circ \mathrm{L}_{\Fanin{}} \circ a_{- 2}\).

\subsubsection{Key Statistics}

Several statistics computed over the dataset are used across all
methods:

\begin{definition}

The \textbf{activation covariance} at layer \(l\) is:
\[A_{l} \coloneqq \EE_{(x,y) \sim \D}\left\lbrack a_{l}(x) \times_{1}\left( a_{l}(x) \right)^{\top} \right\rbrack \approx \frac{1}{n}{\mathbf{a}_{l}}^{\top} \times_{n}\mathbf{a}_{l} \;\;\in \left( C_{l},C_{l} \right)\]

\end{definition}

\begin{definition}

The \textbf{gradient-activation cross-covariance} with activation at
layer \(l\) with $l \in \{-1, -2\}$ is:
\[B_{l} \coloneqq \EE_{(x,y) \sim \D}\left\lbrack \left( - {\nabla_{s(x)}\ell(f(x))} \right) \times_{1}\left( a_{l}(x) \right)^{\top} \right\rbrack \approx - \frac{1}{n}\Gradpreact^{\top} \times_{n}\mathbf{a}_{l} \;\;\in \left( C,C_{l} \right)\]

\end{definition}

\begin{definition}

The \textbf{gradient covariance} is:
\[\covgrad \coloneqq \EE_{(x,y) \sim \D}\left\lbrack \left( \frac{\partial f}{\partial s_{}}(x) \right)^{\top} \times_{C_{L}}\frac{\partial f}{\partial s_{}}(x) \right\rbrack\;\; \in (C,C)\]

\end{definition}

\subsubsection{Residual Gradient (Bottleneck)}

To avoid redundancy between added neurons and existing ones, several
methods project the gradient orthogonally to the space spanned by
the activity of existing neurons. The idea is to find the optimal update
\({\mathrm{d}W}^{\ast}\) to the existing weights \(W\) and then consider
only the residual gradient that cannot be addressed by this update.



\begin{definition}

The \textbf{residual gradient} (or \textbf{bottleneck}) is the residual
after subtracting the contribution of the optimal weight update
\({\mathrm{d}W}^{\ast}\):
\[\left( {\nabla_{s_{}}\loss(f)} \right)^{\bot}(x) \coloneqq - {\nabla_{s_{}}\loss(f)}(x) - \mathrm{L}_{{\mathrm{d}W}^{\ast}}(a_{- 1}(x)) \in (C)\]
In batched form:
\[{\Bottleneck} \coloneqq - \Gradpreact - \mathbf{a}_{- 1} \times_{C_{- 1}}A_{+}^{- 1} \times_{C_{- 1}}{B_{- 1}}^{\top} \in (n,C)\]

\end{definition}

\begin{lemma}\label{lemma:optimal_dW}

The optimal weight update \({\mathrm{d}W}^{\ast}\) that minimizes the
residual gradient is:
\[{\mathrm{d}W}^{\ast} \coloneqq \argmin_{\mathrm{d}W}\left\| {- {\nabla_{s_{}}\loss(f)} - \mathrm{L}_{\mathrm{d}W} \circ a_{- 1}} \right\|_{2} = B_{- 1} \times_{C_{- 1}}{A_{- 1}}^{+}\]

\end{lemma}

\begin{proof}

We solve it for a batch of data points:
\({\mathrm{d}W}^{\ast} = \argmin_{\mathrm{d}W}\left\| {- \Gradpreact - \mathbf{a}_{- 1} \times_{C_{- 1}}{\mathrm{d}W}^{\top}} \right\|_{2}\)
This is a least-squares problem:
\[{{\mathrm{d}W}^{\ast}}^{\top} = \left( {\mathbf{a}_{- 1}}^{\top} \times_{n}\mathbf{a}_{- 1} \right)^{+} \times_{C_{- 1}}{\mathbf{a}_{- 1}}^{\top} \times_{n}\left( - \Gradpreact \right) = {A_{- 1}}^{+} \times_{C_{- 1}}{B_{- 1}}^{\top}\]

\end{proof}

\begin{proposition}\label{prop:bottleneck_orthogonality}

The residual gradient \({\Bottleneck}\) is orthogonal to existing
activations:
\[{\mathbf{a}_{- 1}}^{\top} \times_{n}{\Bottleneck} = 0 \in \left( C_{- 1},C \right)\]

\end{proposition}

\begin{proof}

By construction, \(\left( {\nabla_{s_{}}\loss(f)} \right)^{\bot}\) is
the residual of the least-squares projection of
\(- {\nabla_{s_{}}\loss(f)}\) onto the space spanned by \(a_{- 1}\). The
residual is orthogonal to the projection space.

\end{proof}

\subsection{Existing Methods}

This section presents each method with its motivation and derives its
simplified objective function.

\subsubsection{TINY}

\noindent\textbf{Motivation.} The TINY~\citep{verbockhaven_growing_2024} method seeks to approximate the desired pre-activation update \(-\nabla_{s_{}}\loss(f)\) by changes to existing weights and the contribution of new neurons.
It first computes the best update \(\mathrm{d}W^*\) using existing neurons alone, as in \autoref{lemma:optimal_dW}.
Keeping this update fixed, it then chooses new neurons to approximate the remaining residual \(\Bottleneck\).
This is a sequential least-squares construction; it need not solve the joint optimization over existing and new weights~\citep{verbockhaven_growing_2024}.

\noindent\textbf{Objective Formulation.} TINY approximates
\(\delta_{s_{}}(x) \approx \sigma'(0) \cdot \Fanout{} \times_{\cext}\Fanin{} \times_{C_{- 2}}a_{- 2}(x)\)
(linearizing \(\sigma\) around \(0\)). Dropping the constant
\(\sigma'(0)\), the objective becomes:

\begin{proposition}

\textbf{TINY optimization problem (with projection).} The TINY method~\footnote{As stated in  equation (8) of \citep{verbockhaven_growing_2024}.} solves the following problem for the new weights:

\begin{equation}\label{eq:tiny_projection_optimization}
\argmin_{\Fanin{},\Fanout{}}\left\| {{\Bottleneck} - \mathbf{a}_{- 2} \times_{C_{- 2}}\Fanin{}^{\top} \times_{\cext}\Fanout{}^{\top}} \right\|_{2}
\end{equation}
Without projection, replace \({\Bottleneck}\) with \(- \Gradpreact\).

\end{proposition}

\begin{proof}

We maximize the first-order loss decrease. If $a_{-2}^\top G^\bot \neq 0$, the non degeneracy condition is met and by
\autoref{corollary:argmin_norm_to_argmax_scalarp}:
\[\argmin_{\Fanin{},\Fanout{};\left\| \delta_{s_{}} \right\| \leq 1}{\langle\left( {\nabla_{s_{}}\loss(f)} \right)^{\bot},\delta_{s_{}}\rangle} \propto \argmin_{\Fanin{},\Fanout{}}\left\| {- \left( {\nabla_{s_{}}\loss(f)} \right)^{\bot} - \delta_{s_{}}} \right\|\]
Then substituting the linearized \(\delta_{s_{}}\) and using batched
notation gives the result.

If $a_{-2}^\top G^\bot = 0$, in practice $\deltas = 0$ is selected, which is one of the solutions for both problems.


\end{proof}

\subsubsection{GradMax}

\textbf{Motivation.} GradMax~\citep{evci_gradmax_2022} takes a different
approach: initialize new neurons with zero contribution
(\(\delta_{s_{}} = 0\) by setting \(\Fanin{} = 0\)), then maximize the
gradient magnitude that will be available for subsequent training.

\textbf{Objective Formulation.} With \(\Fanin{} = 0\) and \(\sigma(0) = 0\), we have
\(\postactext = 0\), so \(\nabla_{\Fanout{}}\loss(f) = 0\). The loss
decreases after one gradient step on
\(\left( \Fanin{},\Fanout{} \right)\) is:
\[\loss(f_{\Fanin{} + {\mathrm{d}\Fanin{}}, \Fanout{} + {\mathrm{d}\Fanout{}}}) \approx \loss(f) - {\left\| {\nabla_{\Fanin{}}\loss(f)} \right\|_{2}}^{2} - {\left\| {\nabla_{\Fanout{}}\loss(f)} \right\|_{2}}^{2}\]
As \(\nabla_{\Fanout{}}\loss(f) = 0\) at initialization, GradMax
maximizes \(\left\| {\nabla_{\Fanin{}}\loss(f)} \right\|_{2}\).

\begin{proposition}

\textbf{GradMax optimization problem.} GradMax solves~\footnote{As stated in equation (11) of \citep{evci_gradmax_2022}}:
\[\Fanout{}^{\ast} = \argmax_{{\left\| \Fanout{} \right\|_{2}}^{2} \leq \cext}\left\| {\nabla_{\Fanin{}}\loss(f)} \right\|_{2}\]
such that \(\Fanout{}^{\top} \times_{\cext}\Fanout{} = I_{\cext}\).

\end{proposition}

\noindent\textbf{Simplified Objective}

\begin{lemma}

With \(\Fanin{} = 0\) and \(\sigma'(0) = 1\):
\[\nabla_{\Fanin{}}\loss(f) = \Fanout{}^{\top} \times_{C}\EE_{(x,y) \sim \D}\left\lbrack {\nabla_{s_{}}\ell(f(x))} \times_{1}{a_{- 2}(x)}^{\top} \right\rbrack = - \Fanout{}^{\top} \times_{C}B_{- 2}\]

\end{lemma}

\begin{proposition}\label{prop:gradmax_reduced}

\textbf{GradMax reduced objective.} The GradMax optimization reduces to:

\begin{equation}\label{eq:gradmax_reduced}
\mathcal{J}_{\text{GradMax}}\left( \Fanout{} \right) \coloneqq {\left\| {{B_{- 2}}^{\top} \times_{C}\Fanout{}} \right\|_{2}}^{2}
\end{equation}
such that \(\Fanout{}^{\top}\Fanout{} = I_{\cext}\)~\footnote{This constraint is not explicitly mentioned in the paper but is implicit knowing the objective is optimized through SVD.}. The optimal \(\Fanout{}^{\ast}\) are the leading left-singular vectors
of \(B_{- 2}\).

\end{proposition}

\noindent\textbf{GradMax-Opt.} GradMax-Opt is similar, but it does not include the orthogonality constraint. Hence, instead of using an SVD, it uses projected gradient descent. This allows an alternative initialization where the roles of \(\Fanin{}\) and \(\Fanout{}\) are switched: setting \(\Fanout{} = 0\) and maximizing \(\left\| {\nabla_{\Fanout{}}\loss(f)} \right\|_{2}\).

\subsubsection{SENN}\label{sec:senn}

\textbf{Motivation}

SENN~\citep{mitchell_self-expanding_2024} extends GradMax
by considering \textbf{natural gradient} descent instead of standard
gradient descent. It also avoids redundancy by maximizing the
\emph{increase} in natural gradient norm after neuron addition. They
choose to fix \(\Fanout{} = 0\) at initialization to ensure
\(\delta_{s_{}} = 0\) (unlike GradMax that fixes \(\Fanin{} = 0\)).

The natural gradient step is \(\begin{smallpmatrix}
{\mathrm{d}\Fanout{}} \\
{\mathrm{d}W}
\end{smallpmatrix} \propto - F^{- 1}\nabla_{\begin{smallpmatrix}
\Fanout{} \\
W
\end{smallpmatrix}}\loss(f)\) where \(F\) is the Fisher information matrix
with respect to \(\begin{smallpmatrix}
\Fanout{} \\
W
\end{smallpmatrix}\), the first-order equation becomes (also because
\(\nabla_{\Fanin{}}\loss(f) = 0\)):
\[\loss(f_{W + {\mathrm{d}W},\Fanout{} + {\mathrm{d}\Fanout{}}}) \approx \loss(f) - {\langle\nabla_{\begin{smallpmatrix}
\Fanout{} \\
W
\end{smallpmatrix}}\loss(f),F^{- 1}\nabla_{\begin{smallpmatrix}
\Fanout{} \\
W
\end{smallpmatrix}}\loss(f)\rangle}_{2}\] Hence they want to maximize
\(\eta_{\text{SENN}}(\begin{smallpmatrix}
\Fanout{} \\
W
\end{smallpmatrix}) \coloneqq {\langle\nabla_{\begin{smallpmatrix}
\Fanout{} \\
W
\end{smallpmatrix}}\loss(f),F^{- 1}\nabla_{\begin{smallpmatrix}
\Fanout{} \\
W
\end{smallpmatrix}}\loss(f)\rangle}_{2}\).

\noindent\textbf{Objective Formulation}
Using the KFAC~\citep{martens_kfac_2015,dangelKroneckerfactoredApproximateCurvature2025} approximation
\(F \approx \covgrad \times A_{- 1}\) for the Fisher information matrix:

\begin{definition}

The \textbf{SENN criterion} measures the natural gradient norm:
\[\begin{aligned}
\eta_{\text{SENN}}(\begin{smallpmatrix}
\Fanout{} \\
W
\end{smallpmatrix}) & \coloneqq {\langle\nabla_{\begin{smallpmatrix}
\Fanout{} \\
W
\end{smallpmatrix}}\loss(f),F^{- 1}\nabla_{\begin{smallpmatrix}
\Fanout{} \\
W
\end{smallpmatrix}}\loss(f)\rangle}_{2} \\
\text{Using KFAC:}\quad & = {\langle\nabla_{\begin{smallpmatrix}
\Fanout{} \\
W
\end{smallpmatrix}}\loss(f)\left( A_{- 1}^{\text{conc}} \right)^{- 1},\covgrad^{- 1}\nabla_{\begin{smallpmatrix}
\Fanout{} \\
W
\end{smallpmatrix}}\loss(f)\rangle}_{2}
\end{aligned}\] where:

\begin{itemize}
\item
  \(A_{- 1}^{\text{conc}}\) is the covariance of concatenated
  activations \(\begin{smallpmatrix}
  \postactext \\
  a_{- 1}
  \end{smallpmatrix}\)
\end{itemize}

\end{definition}

SENN maximizes the increase
\(\Delta\eta_{\text{SENN}} \coloneqq \eta_{\text{SENN}}(\begin{smallpmatrix}
\Fanout{} \\
W
\end{smallpmatrix}) - \eta_{\text{SENN}}(W)\).

\noindent\textbf{Simplified Objective}
SENN uses a tractable lower bound
\(\Delta\eta_{\text{SENN}}' \leq \Delta\eta_{\text{SENN}}\) involving
the residual gradient.

\begin{definition}
The \textbf{extended activation covariance} (using exact activations):
\[A_{- 1}^{\text{ext}} \coloneqq \EE_{(x,y) \sim \D}\left\lbrack \postactext(x) \times_{1}{\postactext(x)}^{\top} \right\rbrack \in \left( \cext,\cext \right)\]
\end{definition}

\begin{definition}
The \textbf{residual gradient-activation cross-covariance} (using exact
activations):
\[\partial\Fanout{} \coloneqq \EE_{(x,y) \sim \D}\left\lbrack \left( \nabla_{s(x)}\ell(f(x)) \right)^{\bot} \times_{1}{\postactext(x)}^{\top} \right\rbrack \in \left( C,\cext \right)\]
\end{definition}

\begin{proposition}\label{prop:senn_reduced}
\textbf{SENN reduced objective.~\footnote{As stated in Theorem 3.2 of~\citep{mitchell_self-expanding_2024}.}} The SENN lower bound objective is
\(\eta_{\text{SENN}}(\Fanout{})\) but computing the gradient of
\(\Fanout{}\) backpropagating only the residual gradient (hence the
definition of \(\partial\Fanout{}\) the pseudo gradient of
\(\Fanout{}\)):
\begin{equation}\label{eq:senn_reduced}
\begin{aligned}
\mathcal{J}_{\text{SENN}}\left( \Fanin{} \right) & \coloneqq {\langle\partial\Fanout{} \times_{\cext}{A_{- 1}^{\text{ext}}}^{+},\covgrad^{- 1} \times_{C}\partial\Fanout{}\rangle}_{2} \\
 & = \left\| {\isqrtS \times_{C}\partial\Fanout{} \times_{\cext}\isqrtAext} \right\|_{2}^{2}
\end{aligned}
\end{equation}
\end{proposition}

\subsubsection{NeST}

\noindent\textbf{Motivation}

NeST~\citep{daiNeSTNeuralNetwork2019} considers sparse networks where
connections are added incrementally. NeST provides criteria for both
\textbf{connection activation} (adding a connection in an existing
layer) and \textbf{neuron addition} (adding a new neuron with a single
connection pair).

\textbf{Connection Activation}

NeST considers sparse layers where \(W\) is sparse. To choose which
connection to add, they look at which connection between neuron \(i\) in
layer \(l - 1\) and neuron \(j\) in layer \(l\) would maximally decrease
the loss if added and optimized with a gradient step:
\[\loss(f_{W + {\mathrm{d}W}\lbrack i,j\rbrack}) = \loss(f) + {\langle\nabla_{W\lbrack i,j\rbrack}\loss(f),{\mathrm{d}W}\lbrack i,j\rbrack\rangle} + o\left( \left\| {{\mathrm{d}W}\lbrack i,j\rbrack} \right\|_{2} \right)\]
When choosing
\({\mathrm{d}W}\lbrack i,j\rbrack \propto - \nabla_{W\lbrack i,j\rbrack}\loss(f)\),
this becomes:
\[\loss(f_{W + {\mathrm{d}W}\lbrack i,j\rbrack}) = \loss(f) - {\left\| {\nabla_{W\lbrack i,j\rbrack}\loss(f)} \right\|_{2}}^{2} + o\left( \left\| {{\mathrm{d}W}\lbrack i,j\rbrack} \right\|_{2} \right)\]
Note that NeST justifies this choice with Hebbian theory and considers
\(\left| {\nabla_{W\lbrack i,j\rbrack}\loss(f)} \right|\).

\begin{proposition}\label{prop:nest_connection}

\textbf{NeST connection activation.} For adding a connection in an
existing layer, NeST selects:
\[i^{\ast},j^{\ast} = \argmax_{i,j}\left| B_{- 1} \right|\lbrack j,i\rbrack\]

\end{proposition}

\noindent\textbf{Neuron Addition}
When adding a new neuron, NeST only adds a few connections between pairs
of neurons \((i,j)\) with neuron \(i\) in layer \(l - 2\) and neuron
\(j\) in layer \(l\). They set \(\Fanin{}\) and \(\Fanout{}\) to zero
except for the values at indices \((i)\) and \((j)\) respectively. This
gives:
\[\delta_{s_{}}(x)\lbrack j\rbrack = \Fanout{}\lbrack j\rbrack\sigma(\Fanin{}\lbrack i\rbrack a_{- 2}(x)\lbrack i\rbrack) \approx \Fanout{}\lbrack j\rbrack\sigma'(0)\Fanin{}\lbrack i\rbrack a_{- 2}(x)\lbrack i\rbrack\]
Dropping the constant \(\sigma'(0)\), they optimize the decrease in loss
after the addition: \[\begin{aligned}
\loss(s_{} + \delta_{s_{}}) & = \loss(s_{}) + {\langle\nabla_{s_{}\lbrack j\rbrack}\loss(s_{}),\delta_{s_{}}\lbrack j\rbrack\rangle} + o\left( \left\| {\delta_{s_{}}\lbrack j\rbrack} \right\| \right) \\
 & \approx \loss(s_{}) + \nabla_{s_{}\lbrack j\rbrack}\loss(s_{}) \cdot \Fanout{}\lbrack j\rbrack \cdot \Fanin{}\lbrack i\rbrack \cdot a_{- 2}\lbrack i\rbrack
\end{aligned}\]
Hence they minimize
\(\nabla_{s_{}\lbrack j\rbrack}\loss(s_{}) \cdot a_{- 2}\lbrack i\rbrack \cdot \Fanout{}\lbrack j\rbrack \cdot \Fanin{}\lbrack i\rbrack\)
under the constraint that
\(\left| {\Fanout{}\lbrack j\rbrack \cdot \Fanin{}\lbrack i\rbrack} \right|\)
is small.

\begin{proposition}\label{prop:nest_objective}

\textbf{NeST neuron addition.} For adding a new neuron, NeST selects:
\begin{equation}\label{eq:nest_neuron_metric}
i^{\ast},j^{\ast} = \argmax_{i,j}\left| B_{- 2} \right|\lbrack j,i\rbrack
\end{equation}
to add a single connection, or the pairs in the top \(\beta\) fraction
to add several connections.

\end{proposition}

\noindent\textbf{Initialization} Once \(\left( i^{\ast},j^{\ast} \right)\) is chosen, NeST initializes the weights so that the change \(\delta_{s_{}}\lbrack j\rbrack\) is similar to that created by a gradient step on a virtual connection
between \(a_{- 2}\lbrack i\rbrack\) and \(s_{}\lbrack j\rbrack\), yielding a loss decrease of \({\left\| {B_{- 2}\left\lbrack j^{\ast},i^{\ast} \right\rbrack} \right\|_{2}}^{2}\). The sign is randomized:

\begin{proposition}\label{prop:nest_init}

\textbf{NeST initialization.} Once \(\left( i^{\ast},j^{\ast} \right)\)
is selected, NeST initializes: \[\begin{aligned}
\varepsilon & \sim \text{ Rademacher }\text{i.e.}\text{Uniform}\left( \left\{ - 1,1 \right\} \right) \\
\Fanin{}\left\lbrack i^{\ast} \right\rbrack & = \varepsilon \cdot \text{ sign}\left( B_{- 2}\left\lbrack j^{\ast},i^{\ast} \right\rbrack \right) \cdot \sqrt{\left| {B_{- 2}\left\lbrack j^{\ast},i^{\ast} \right\rbrack} \right|} \\
\Fanout{}\left\lbrack j^{\ast} \right\rbrack & = \varepsilon \cdot \sqrt{\left| {B_{- 2}\left\lbrack j^{\ast},i^{\ast} \right\rbrack} \right|}
\end{aligned}\]
All other entries of \(\Fanin{}\) and \(\Fanout{}\) are zero.

\textbf{Note:} The NeST paper seems to align the product
\(\Fanout{}\left\lbrack j^{\ast} \right\rbrack \cdot \Fanin{}\left\lbrack i^{\ast} \right\rbrack\)
with \(- B_{- 2}\left\lbrack j^{\ast},i^{\ast} \right\rbrack\) instead
of \(B_{- 2}\left\lbrack j^{\ast},i^{\ast} \right\rbrack\), which would
maximize the loss instead of minimizing it. We present here what we
believe is the intended version.

\end{proposition}

\subsubsection{NORTH}

\textbf{Motivation}

NORTH~\citep{maile_when_2022} takes a geometric perspective:
new neurons should add \textbf{new directions} in activation space,
orthogonal to existing activations.

\textbf{Objective Formulation}

NORTH seeks \(\fanin{}\) such that the new activations
\(\Postactext = \sigma(\mathbf{a}_{- 2}\fanin{}^{\top})\) are orthogonal
to existing ones, which is equivalent to lying in
\(\ker({\mathbf{a}_{- 1}}^{\top})\).

\begin{proposition}\label{prop:north_construction}

\textbf{NORTH-Pre construction.} Let
\(V \in \left( n,n - C_{- 1} \right)\) be an orthonormal basis of
\(\ker({\mathbf{s}_{- 1}}^{\top})\) (computed via SVD). For random
directions \(R \in \left( n - C_{- 1},\cext \right)\):

\begin{equation}\label{eq:north_proposition}
\Fanin{}^{\top} = A_{- 2}^{- 1} \times_{C_{- 2}}{\mathbf{a}_{- 2}}^{\top} \times_{n}V \times_{n - C_{- 1}}R
\end{equation}

\end{proposition}

\paragraph{Justification}

They use the approximation \(\sigma \approx \text{ Identity}\) to relax
the problem to replace \(\Postactext\) with
\(\Preactext = \mathbf{a}_{- 2} \times_{C_{- 2}}\fanin{}^{\top}\) and
\(\ker({\mathbf{a}_{- 1}}^{\top})\) with
\(\ker({\mathbf{s}_{- 1}}^{\top})\).

First they compute a basis \(V \in \left( n,n - C_{- 1} \right)\) of
\(\ker({\mathbf{s}_{- 1}}^{\top})\) using SVD. Then they generate random
combinations of the basis vectors \(V \times_{n - C_{- 1}}r\) for random
\(r\) as candidates for \(\Preactext\). For each candidate, they solve
for \(\fanin{}\) such that
\(\mathbf{a}_{- 2} \times_{C_{- 2}}\fanin{}^{\top} = V \times_{n - C_{- 1}}r\)
which gives:
\[\fanin{}^{\top} = \left( \mathbf{a}_{- 2} \right)^{+} \times_{n}Vr\]

Or for a batched version with \(\cext\) candidates
\(R \in \left( n - C_{- 1},\cext \right)\) with orthogonal columns:
\[\begin{aligned}
\Fanin{}^{\top} & = \left( \mathbf{a}_{- 2} \right)^{+} \times_{n}V \times_{n - C_{- 1}}R \in \left( C_{- 2},\cext \right) \\
 & = \left( {\mathbf{a}_{- 2}}^{\top} \times_{n}\mathbf{a}_{- 2} \right)^{+} \times_{C_{- 2}}{\mathbf{a}_{- 2}}^{\top} \times_{n}V \times_{n - C_{- 1}}R \\
 & \propto A_{- 2}^{- 1} \times_{C_{- 2}}{\mathbf{a}_{- 2}}^{\top} \times_{n}V \times_{n - C_{- 1}}R
\end{aligned}\] (\(\propto\) because we skip the \(\frac{1}{n}\)
factor.) Then they select the best candidates according to the
orthogonality of
\(\Postactext = \sigma(\mathbf{a}_{- 2}\fanin{}^{\top})\) to
\(\mathbf{a}_{- 1}\) using this time the exact activation function
\(\sigma\).

\subsection{Unified Framework}

For completeness, we now restate \autoref{thm:unified_obj_main_text}, which shows that all methods can be understood as special cases of a
common optimization objective.

\begin{theorem}\label{thm:unified_objective}

\textbf{Unified objective for neuron addition.} All neuron addition
methods studied in this document can be derived from the following
optimization problem:

\begin{subequations}\label{eq:unified_objective}
\begin{align}
\Fanin{}^{\ast},\Fanout{}^{\ast} & = \argmax_{\Fanin{},\Fanout{},\left\| {\sqrtAext \times_{\cext}\Fanout{}^{\top} \times_{C}\sqrtS} \right\|_{2} \leq 1}{\langle T \times_{C}\isqrtS,\frac{1}{n}\Postactext \times_{\cext}\Fanout{}^{\top} \times_{C}\sqrtS\rangle} \label{eq:unified_objective_scalarp} \\
\shortintertext{if we linearize $\sigma$ around $0$ and $\transp{\postact[-2]}T \neq 0$:}
 & \propto \argmin_{\Fanin{},\Fanout{}}\left\| {T \times_{C}\isqrtS - \sigma(\mathbf{a}_{- 2} \times_{C_{- 2}}\Fanin{}^{\top}) \times_{\cext}\Fanout{}^{\top} \times_{C}\sqrtS} \right\|_{2} \label{eq:unified_objective_norm}
\end{align}
\end{subequations}
where:
\begin{itemize}
\item
  \(T \in (n,C)\) is the target (negative gradient \(- \Gradpreact\) or
  residual gradient \({\Bottleneck}\))
\item
  \(\covgrad \in (C,C)\) is the gradient weighting metric
\item
  \(\sigma\) is an activation function
\item
  Additional constraints may be imposed on \(\Fanin{}\) and
  \(\Fanout{}\)
\end{itemize}

Note that if $\transp{\postact[-2]}T = 0$ (under the linearization assumption), $(0, 0)$ is an optimal solution in both cases and is the one selected in practice.
The following table summarizes the specialization for each method:
\begin{longtable}[]{@{}ccccccccc@{}}
\caption{Neuron addition methods as special cases of the unified
objective in \autoref{eq:unified_objective}.}\tabularnewline
\toprule\noalign{}
\textbf{Method} & \(T\) & \(A_{- 1}^{\text{ext}}\) & \(\covgrad\) &
\textbf{Extra constraints} & \(\Fanin{}\) & \(\Fanout{}\) &
\(\postactext \perp a_{- 1}\) & \(\sigma\) \\
\midrule\noalign{}
\endfirsthead
\toprule\noalign{}
\textbf{Method} & \(T\) & \(A_{- 1}^{\text{ext}}\) & \(\covgrad\) &
\textbf{Extra constraints} & \(\Fanin{}\) & \(\Fanout{}\) &
\(\postactext \perp a_{- 1}\) & \(\sigma\) \\
\midrule\noalign{}
\endhead
\bottomrule\noalign{}
\endlastfoot
Full & \({\Bottleneck}\) & \(A_{- 1}^{\text{ext}}\) & \(\covgrad\) &
\(\varnothing\) & \(\Fanin{}^{\ast}\) & \(\Fanout{}^{\ast}\) &
\(\checkmark\) & Any \\
TINY (proj.) & \({\Bottleneck}\) & \(\Fanin{}A_{- 2}\Fanin{}^{\top}\) &
\(I_{C}\) & \(\varnothing\) & \(\Fanin{}^{\ast}\) & \(\Fanout{}^{\ast}\)
& \(\checkmark\) & Identity \\
TINY (no proj.) & \(- \Gradpreact\) & \(\Fanin{}A_{- 2}\Fanin{}^{\top}\)
& \(I_{C}\) & \(\varnothing\) & \(\Fanin{}^{\ast}\) &
\(\Fanout{}^{\ast}\) &
\(\times\) & Identity \\
GradMax & \(- \Gradpreact\) & \(\Fanin{}\Fanin{}^{\top}\) & \(I_{C}\) &
\(\Fanin{}\Fanin{}^{\top} = I_{\cext};\Fanout{}^{\top}\Fanout{} = I_{\cext}\)
& \(0\) & \(\Fanout{}^{\ast}\) &
\(\times\) & Any \\
SENN & \({\Bottleneck}\) & \(A_{- 1}^{\text{ext}}\) & \(\covgrad\) &
\(\varnothing\) & \(\Fanin{}^{\ast}\) & \(0\) &
\(\checkmark\) & Any \\
NeST & \(- \Gradpreact\) & \(\Fanin{}\Fanin{}^{\top}\) & \(I_{C}\) &
\(\left\| \Fanin{} \right\|_{0} = \left\| \Fanout{} \right\|_{0} = 1\) &
\(\Fanin{}^{\ast}\) & \(\Fanout{}^{\ast}\) &
\(\times\) & Identity \\
\end{longtable}

\phantomsection\label{tab:comparison}

\noindent Notes:

\begin{itemize}
\item
  Assuming \(A_{- 1}^{\text{ext }} = \Fanin{}A_{- 2}\Fanin{}^{\top}\) in
  \autoref{eq:unified_objective_scalarp} is equivalent to linearizing
  \(\sigma\) around \(0\) in \autoref{eq:unified_objective_norm}.
\item
  The GradMax assumption
  \(A_{- 1}^{\text{ext }} = \Fanin{}\Fanin{}^{\top}\) is like assuming
  the input data is whitened in addition to linearizing \(\sigma\)
  around \(0\).
\item
  The equivalence between the two optimization problems is correct
  either if we use the first-order approximation of \(\sigma\) or if we
  allow to have different proportionality constants for \(\Fanin{}\) and
  \(\Fanout{}\).
\end{itemize}

\end{theorem}

\begin{proof}

The derivations for each method are provided in the following sections.
We only prove the equivalence between the minimization and maximization
forms. By \autoref{lemma:argmin_norm_to_argmax_scalarp} from the
preliminaries, we have:
\begin{equation}
\begin{aligned}
\Fanin{}^{\ast},\Fanout{}^{\ast} & = \argmin_{\Fanin{},\Fanout{}}\left\| {T \times_{C}\isqrtS - \sigma(\mathbf{a}_{- 2} \times_{C_{- 2}}\Fanin{}^{\top}) \times_{\cext}\Fanout{}^{\top} \times_{C}\sqrtS} \right\|_{2} \\
 & = \argmin_{\Fanin{},\Fanout{}}\frac{1}{\sqrt{n}}\left\| {T \times_{C}\isqrtS - \Postactext \times_{\cext}\Fanout{}^{\top} \times_{C}\sqrtS} \right\|_{2}
\end{aligned}
\end{equation}
Linearizing $\sigma$ around $0$, $\Postactext = \mathbf{a}_{- 2} \times_{C_{- 2}}\Fanin{}^{\top}$ is linear in $\Fanin{}$, so $f:\left( \Fanin{},\Fanout{} \right) \mapsto \frac{1}{\sqrt{n}}\Postactext\Fanout{}^{\top}\sqrtS$ is homogeneous of degree $2$, $n:\left( \Fanin{},\Fanout{} \right) \mapsto \left( \Fanin{}, - \Fanout{} \right)$ is a valid sign-flip operator. As we assume $\transp{\postact[-2]}T \neq 0$ the non-degeneracy condition is met and \autoref{corollary:argmin_norm_to_argmax_scalarp} applies:
\begin{equation}
\begin{aligned}
 & \propto \argmax_{\Fanin{},\Fanout{},\left\| {\frac{1}{\sqrt{n}}\Postactext \times_{\cext}\Fanout{}^{\top} \times_{C}\sqrtS} \right\|_{2} \leq 1}{\langle\frac{1}{\sqrt{n}}T \times_{C}\isqrtS,\frac{1}{\sqrt{n}}\Postactext \times_{\cext}\Fanout{}^{\top} \times_{C}\sqrtS\rangle} \\
 & = \argmax_{\Fanin{},\Fanout{},\left\| {\sqrtAext \times_{\cext}\Fanout{}^{\top} \times_{C}\sqrtS} \right\|_{2} \leq 1}{\langle T \times_{C}\isqrtS,\frac{1}{n}\Postactext \times_{\cext}\Fanout{}^{\top} \times_{C}\sqrtS\rangle}
\end{aligned}
\end{equation}

the last step using \autoref{theorem:rimannian_norm} together with
\(A_{- 1}^{\text{ext }} = \frac{1}{n}\Postactext^{\top} \times_{n}\Postactext\).
This equivalence is only needed for the norm-form methods, i.e., TINY,
which assumes \(\sigma \approx \text{ Identity}\); GradMax, NeST and
SENN are obtained directly from the maximization form
\autoref{eq:unified_objective_scalarp}.

\end{proof}

\subsubsection{Derivation of Each Method}

We now prove that each method can be derived from the unified objective
\autoref{eq:unified_objective} by appropriate specialization.

\noindent\textbf{Derivation of TINY}

\begin{proposition}\label{prop:tiny_derivation}

\textbf{TINY from unified objective.} Setting \(\covgrad = I_{C}\) and
\(\sigma \approx \text{ Identity}\) in \autoref{eq:unified_objective_norm}
yields the TINY optimization problem.

\end{proposition}

\begin{proof}

With \(\covgrad = I_{C}\), we have \(\isqrtS = I_{C}\). The unified
objective becomes:
\[\argmin_{\Fanin{},\Fanout{}}\left\| {T - \sigma(\mathbf{a}_{- 2} \times_{C_{- 2}}\Fanin{}^{\top}) \times_{\cext}\Fanout{}^{\top}} \right\|_{2}\]
Using the linear approximation \(\sigma \approx \text{ Identity}\)
(which is equivalent to
\(A_{- 1}^{\text{ext }} = \Fanin{}A_{- 2}\Fanin{}^{\top}\)), this
becomes:
\[\argmin_{\Fanin{},\Fanout{}}\left\| {T - \mathbf{a}_{- 2} \times_{C_{- 2}}\Fanin{}^{\top} \times_{\cext}\Fanout{}^{\top}} \right\|_{2}\]
With \(T = {\Bottleneck}\), this is exactly
\autoref{eq:tiny_projection_optimization}. With \(T = - \Gradpreact\), this
is TINY without projection.

\end{proof}

\noindent\textbf{Derivation of GradMax}

\begin{proposition}\label{prop:gradmax_derivation}

\textbf{GradMax from unified objective.} Setting
\(T \leftarrow - \Gradpreact\), \(A_{- 2} \leftarrow I_{C_{- 2}}\),
\(\covgrad \leftarrow I_{C}\), with the constraint
\(\Fanout{}^{\top}\Fanout{} = I_{\cext}\), and then setting
\(\Fanin{} = 0\) after optimization yields GradMax.
\end{proposition}
This proof is largely inspired by the work of~\citep{verbockhaven_growing_2024}.

\begin{proof}

Starting from \autoref{eq:unified_objective_scalarp} with
\(\covgrad \leftarrow I_{C}\) and \(T \leftarrow - \Gradpreact\), the
objective is: \[\begin{aligned}
\argmax_{\Fanin{},\Fanout{},\left\| {\left( \Fanin{}\Fanin{}^{\top} \right)^{\frac{1}{2}}\Fanout{}^{\top}} \right\|_{2} \leq 1} & {\langle - \Gradpreact,\Postactext \times_{\cext}\Fanout{}^{\top}\rangle}
\end{aligned}\] Using \autoref{theorem:rimannian_norm} with
\(A \leftarrow \Fanin{}^{\top}\) we have that
\(\left\| {\left( \Fanin{}\Fanin{}^{\top} \right)^{\frac{1}{2}}\Fanout{}^{\top}} \right\|_{2} = \left\| {\Fanin{}^{\top}\Fanout{}^{\top}} \right\|_{2}\).
In addition, we replace \(\Postactext\) with its linear approximation
\(\mathbf{a}_{- 2}\Fanin{}^{\top}\). However, this approximation is not a
strong assumption, as it is only used during backward at
\(\sigma(\mathbf{a}_{- 2}\Fanin{}^{\top}) = \sigma(0)\) as GradMax set
\(\Fanin{} = 0\). Therefore as long as \(\sigma'(0) = 1\) the
approximation is exact. We get:

\[\begin{aligned}
 = \argmax_{\left\| {\Fanin{}^{\top} \times_{\cext}\Fanout{}^{\top}} \right\|_{2} \leq 1} & {\langle - \Gradpreact,\frac{1}{n}\mathbf{a}_{- 2}\Fanin{}^{\top} \times_{\cext}\Fanout{}^{\top}\rangle} \\
 = \argmax_{\left\| {\Fanin{}^{\top} \times_{\cext}\Fanout{}^{\top}} \right\|_{2} \leq 1} & {\langle{\frac{1}{n}\mathbf{a}_{- 2}}^{\top} \times_{n}\left( - \Gradpreact \right),\Fanin{}^{\top} \times_{\cext}\Fanout{}^{\top}\rangle} \\
 = \argmax_{\left\| {\Fanin{}^{\top} \times_{\cext}\Fanout{}^{\top}} \right\|_{2} \leq 1} & {\langle{B_{- 2}}^{\top},\Fanin{}^{\top} \times_{\cext}\Fanout{}^{\top}\rangle} \\
\text{Using}\ {}: \propto \argmin_{\Fanin{},\Fanout{}} & \left\| {{B_{- 2}}^{\top} - \Fanin{}^{\top} \times_{\cext}\Fanout{}^{\top}} \right\|_{2}
\end{aligned}\]
This is equivalent to a low-rank approximation of \({B_{- 2}}^{\top}\).
By \autoref{theorem:separable_optimization}, optimizing over \(\Fanout{}\)
with
\(\Fanin{}^{\top} \leftarrow {\Fanin{}^{\ast}}^{\top} = {B_{- 2}}^{\top}\Fanout{}\left( \Fanout{}^{\top}\Fanout{} \right)^{+} = {B_{- 2}}^{\top}\Fanout{}\)
gives: \[\begin{aligned}
 & \propto \argmax_{{\left\| \Fanout{} \right\|_{2}}^{2} \leq \cext}\left\| {{B_{- 2}}^{\top}\Fanout{}\Fanout{}^{\top}} \right\|_{2} \\
\text{Using}\ {\Fanout{}^{\top}\Fanout{} = I_{\cext}}:\qquad & = \argmax_{{\left\| \Fanout{} \right\|_{2}}^{2} \leq \cext}\left\| {{B_{- 2}}^{\top}\Fanout{}} \right\|_{2}
\end{aligned}\]
This is exactly the GradMax objective
\(\mathcal{J}_{\text{GradMax}}\left( \Fanout{} \right)\) from
\autoref{eq:gradmax_reduced}. Finally, setting \(\Fanin{} = 0\) after
optimization matches GradMax's initialization.

\end{proof}

\noindent\textbf{Derivation of SENN}

\begin{proposition}\label{prop:senn_derivation}

\textbf{SENN from unified objective.} The SENN objective
\(\mathcal{J}_{\text{SENN}}\) is recovered from the maximization form
\autoref{eq:unified_objective_scalarp} with \(T = {\Bottleneck}\) and the
full gradient covariance \(\covgrad\), setting \(\Fanout{} = 0\) after
optimization.

\end{proposition}

\begin{proof}

We start from the maximization form \autoref{eq:unified_objective_scalarp}
with \(T = {\Bottleneck}\) and the full metric \(\covgrad\). Since
\({\langle{\Bottleneck} \times_{C}\isqrtS,X \times_{C}\sqrtS\rangle} = {\langle{\Bottleneck},X\rangle}\)
for any \(X\), the \(\covgrad\) factors cancel in the objective (the
metric only enters through the constraint), which becomes
\[\argmax_{\Fanin{},\Fanout{},\left\| {\sqrtAext \times_{\cext}\Fanout{}^{\top} \times_{C}\sqrtS} \right\|_{2} \leq 1}\frac{1}{n}{\langle{\Bottleneck},\Postactext \times_{\cext}\Fanout{}^{\top}\rangle}.\]
Fix \(\Fanin{}\) and let
\(M \coloneqq \sqrtAext \times_{\cext}\Fanout{}^{\top} \times_{C}\sqrtS\),
so the constraint reads \(\left\| M \right\|_{2} \leq 1\) and
\(\Fanout{}^{\top} = \isqrtAext \times_{\cext}M \times_{C}\isqrtS\).
Moving \(\Postactext\) across the inner product brings out the
cross-covariance
\(\partial\Fanout{} = \frac{1}{n}{\Bottleneck}^{\top} \times_{n}\Postactext\),
and since \(\isqrtAext\) and \(\isqrtS\) are symmetric,
\[\begin{aligned}
\frac{1}{n}{\langle{\Bottleneck},\Postactext \times_{\cext}\Fanout{}^{\top}\rangle} & = {\langle\frac{1}{n}\Postactext^{\top} \times_{n}{\Bottleneck},\Fanout{}^{\top}\rangle} = {\langle{\partial\Fanout{}}^{\top},\Fanout{}^{\top}\rangle} \\
 & = {\langle{\partial\Fanout{}}^{\top},\isqrtAext \times_{\cext}M \times_{C}\isqrtS\rangle} \\
 & = {\langle\isqrtAext \times_{\cext}{\partial\Fanout{}}^{\top} \times_{C}\isqrtS,M\rangle},
\end{aligned}\] whose maximum over \(\left\| M \right\|_{2} \leq 1\) is,
by Cauchy--Schwarz (attained at
\(M \propto \isqrtAext{\partial\Fanout{}}^{\top}\isqrtS\)),
\[\left\| {\isqrtAext \times_{\cext}{\partial\Fanout{}}^{\top} \times_{C}\isqrtS} \right\|_{2} = \left\| {\isqrtS \times_{C}\partial\Fanout{} \times_{\cext}\isqrtAext} \right\|_{2} = \sqrt{\mathcal{J}_{\text{SENN}}\left( \Fanin{} \right)}\]
with \(\mathcal{J}_{\text{SENN}}\) as in \autoref{eq:senn_reduced}. The
outer maximization over \(\Fanin{}\) is thus
\(\argmax_{\Fanin{}}\mathcal{J}_{\text{SENN}}\left( \Fanin{} \right)\),
the SENN objective. As in \autoref{prop:senn_reduced} the optimal
\(\Fanout{}\) only serves to score \(\Fanin{}\); SENN keeps
\(\Fanout{} = 0\) at initialization to preserve the function.

\end{proof}

\noindent\textbf{Derivation of NeST}

\begin{proposition}\label{prop:nest_derivation}

\textbf{NeST from unified objective.} Setting
\(T \leftarrow - \Gradpreact\), \(\covgrad \leftarrow I_{C}\),
linearizing \(\sigma\) and adding the sparsity constraint
\(\left\| \Fanin{} \right\|_{0} = \left\| \Fanout{} \right\|_{0} = 1\)
yields NeST.

\end{proposition}

\begin{proof}

Starting from \autoref{eq:unified_objective_scalarp} with
\(A_{- 1}^{\text{ext }} \leftarrow \Fanin{}\Fanin{}^{\top}\),
\(\covgrad \leftarrow I_{C}\), \(T \leftarrow - \Gradpreact\) and
\(\sigma = \text{Id}\), the objective is: \[\begin{aligned}
\argmax_{\Fanin{},\Fanout{},\left\| {\left( \Fanin{}\Fanin{}^{\top} \right)^{\frac{1}{2}}\Fanout{}^{\top}} \right\|_{2} \leq 1,\left\| \Fanin{} \right\|_{0} = \left\| \Fanout{} \right\|_{0} = 1} & {\langle - \Gradpreact,\mathbf{a}_{- 2} \times_{C_{- 2}}\Fanin{}^{\top} \times_{\cext}\Fanout{}^{\top}\rangle} \\
 = \argmax_{\Fanin{},\Fanout{},\left\| {\Fanin{}^{\top}\Fanout{}^{\top}} \right\|_{2} \leq 1,\left\| \Fanin{} \right\|_{0} = \left\| \Fanout{} \right\|_{0} = 1} & {\langle{B_{- 2}}^{\top},\Fanin{}^{\top} \times_{\cext}\Fanout{}^{\top}\rangle}
\end{aligned}\]
To maximize the decrease in the objective, we choose:
\[i^{\ast},j^{\ast} = \argmax_{i,j}\left| {B_{- 2}\lbrack j,i\rbrack} \right|\]
and set:
\(\Fanin{}\left\lbrack i^{\ast} \right\rbrack = \varepsilon{\operatorname{sign}(B_{- 2}\left\lbrack j^{\ast},i^{\ast} \right\rbrack)}\sqrt{\left| {B_{- 2}\left\lbrack j^{\ast},i^{\ast} \right\rbrack} \right|}\)
and
\(\Fanout{}\left\lbrack j^{\ast} \right\rbrack = \varepsilon\sqrt{\left| {B_{- 2}\left\lbrack j^{\ast},i^{\ast} \right\rbrack} \right|}\)
for random \(\varepsilon \in \left\{ - 1,1 \right\}\). This is exactly
the NeST selection criterion from \autoref{prop:nest_objective}. (Note that
we don't respect the exact scaling constraint
\(\left\| {\Fanin{}^{\top}\Fanout{}^{\top}} \right\|_{2} \leq 1\), but
everything is done up to scaling, so this is not an issue.)

\end{proof}

\subsubsection{NORTH Satisfaction}

\begin{proposition}\label{prop:north_satisfaction}

\textbf{Implicit orthogonality in projected methods.} Methods using the
residual gradient \(T = {\Bottleneck}\) implicitly encourage new
activations to be orthogonal to existing ones, similar to NORTH's
explicit construction.

\end{proposition}

\begin{proposition}

\(\fanin{}\) selected by any method with \(T \leftarrow {\Bottleneck}\)
and \(A_{- 1}^{\text{ext }} \leftarrow \Fanin{}A_{- 2}\Fanin{}^{\top}\)
are in the set of \(\fanin{}\) constructed by NORTH-Pre.

\end{proposition}

\begin{proof}

By \autoref{prop:bottleneck_orthogonality},
\({\mathbf{a}_{- 1}}^{\top} \times_{n}{\Bottleneck} = 0\), so
\({\Bottleneck}\) spans directions orthogonal to \(\mathbf{a}_{- 1}\),
similarly to \(V\), which spans directions orthogonal to
\(\mathbf{s}_{- 1}\) (the same under the first-order approximation,
though further from NORTH's actual goal).

For those methods, the optimal solution satisfies:
\[\mathbf{a}_{- 2} \times_{C_{- 2}}\Fanin{}^{\top} \times_{\cext}\Fanout{}^{\top} \approx {\Bottleneck}\]
Hence (\(\propto\) because we skip the \(\frac{1}{n}\) factor):
\[\Fanin{}^{\top} \propto {A_{- 2}}^{+} \times_{C_{- 2}}{\mathbf{a}_{- 2}}^{\top} \times_{n}{\Bottleneck} \times_{C}{\Fanout{}^{\top}}^{+}\]
Comparing with NORTH-Pre's construction in \autoref{eq:north_proposition}:
\[\Fanin{}^{\top} = A_{- 2}^{- 1} \times_{C_{- 2}}{\mathbf{a}_{- 2}}^{\top} \times_{n}V \times_{n - C_{- 1}}R\]
The structures are identical with \(V \leftarrow {\Bottleneck}\) and
\(R \leftarrow {\Fanout{}^{\top}}^{+}\).

\end{proof}


%%%%%%%%%%%%%%%%%%%%%%%%%%%%%%%%%%%%%%%%%%%%%%%%%%%%%%%%%%%%%%%%%%%%%%%%%%%%%%%%%%%%%%%%%%%%%%%%%%%%%%%%%%%%%%%%%%%%%%%%%%%%%%%%%%%%%%%%%%%%%%%%%%%%%%%%%%%%%%%%%%%%%%%%%%%%%%%%%%%%%%%%%%%%%%%%%%%%%%%%%%%%%%%%%%%%%%%%%%%%%%%%%%%%%%%%%%%%%%%%%%%%%%%%%%%%%%%%%%%%%%%%%%%%%%%%%%%%%%%%%%%%%%%%%%%%%%%%%%%%%%%%%%%%%%%%%%%%%%%%%%%%%%%%%%%%%%%%%%%%%%%%%%%%%%%%%%%%%%%%%%%%%%%%%%%%%%%%%%%%%%%%%%%%%%


\end{appendices}

%%===========================================================================================%%
%% If you are submitting to one of the Nature Portfolio journals, using the eJP submission   %%
%% system, please include the references within the manuscript file itself. You may do this  %%
%% by copying the reference list from your .bbl file, paste it into the main manuscript .tex %%
%% file, and delete the associated \verb+\bibliography+ commands.                            %%
%%===========================================================================================%%

\bibliography{bibliography}% common bib file
%% if required, the content of .bbl file can be included here once bbl is generated
%%\input main.bbl

\end{document}
